\documentclass[a4paper,11pt,oneside,openany]{ltjsbook}

% ============================================================
% 日本語・書体
% ============================================================
\usepackage{luatexja-fontspec}
\usepackage{fontspec}
\setmainfont{DejaVu Serif}
\setsansfont{DejaVu Sans}
\setmonofont{DejaVu Sans Mono}
\setmainjfont{Noto Serif CJK JP}
\setsansjfont{Noto Sans CJK JP}
\setmonojfont{Noto Sans Mono CJK JP}

% ============================================================
% 数式・表・図
% ============================================================
\usepackage{amsmath,amssymb,mathtools,bm}
\usepackage{booktabs,longtable,tabularx,array,multirow}
\usepackage{graphicx}
\usepackage{caption,subcaption}
\usepackage{tikz}
\usetikzlibrary{arrows.meta,positioning,shapes.geometric,fit,calc,backgrounds}
\usepackage[most]{tcolorbox}
\usepackage{enumitem}
\usepackage{fancyvrb}
\usepackage{listings}
\usepackage{siunitx}
\usepackage{geometry}
\usepackage{fancyhdr}
\usepackage{titlesec}
\usepackage{hyperref}
\usepackage[nameinlink,noabbrev]{cleveref}
\usepackage{bookmark}
\usepackage{etoolbox}

\geometry{top=23mm,bottom=24mm,left=24mm,right=24mm,headheight=16pt,headsep=8mm,footskip=14mm}
\setlength{\parindent}{1em}
\setlength{\parskip}{0.15em}
\setlist[itemize]{topsep=0.4em,itemsep=0.15em,leftmargin=2.2em}
\setlist[enumerate]{topsep=0.4em,itemsep=0.15em,leftmargin=2.4em}
\renewcommand{\arraystretch}{1.18}
\sisetup{detect-all,group-separator={,},group-minimum-digits=4}

% ============================================================
% 配色・見出し
% ============================================================
\definecolor{DeepNavy}{HTML}{17324D}
\definecolor{Teal}{HTML}{177E89}
\definecolor{Orange}{HTML}{D97706}
\definecolor{Brick}{HTML}{A33A2B}
\definecolor{Forest}{HTML}{2C6E49}
\definecolor{Purple}{HTML}{6B4C9A}
\definecolor{LightBlue}{HTML}{EDF5FA}
\definecolor{LightTeal}{HTML}{EDF8F7}
\definecolor{LightOrange}{HTML}{FFF6E8}
\definecolor{LightRed}{HTML}{FCEFED}
\definecolor{LightGray}{HTML}{F5F6F7}
\definecolor{MidGray}{HTML}{68727D}
\definecolor{DarkGray}{HTML}{2F343A}

\hypersetup{
  colorlinks=true,
  linkcolor=DeepNavy,
  citecolor=Forest,
  urlcolor=Teal,
  pdftitle={カオス同期とTM読み出しによる情報圧縮 数値実験テキスト},
  pdfauthor={長松蒼空},
  pdfsubject={一般化Boole写像・Cauchy分布・TM基底・同期・情報圧縮},
  pdfkeywords={chaotic synchronization, Cauchy distribution, Takenaka-Malmquist basis, Koopman, information compression}
}

\titleformat{\chapter}[display]
  {\normalfont\sffamily\bfseries\color{DeepNavy}}
  {\filleft\Large 第\thechapter 章}
  {0.4ex}
  {\titlerule[1.1pt]\vspace{1.1ex}\Huge}
  [\vspace{0.8ex}\titlerule]
\titleformat{\section}{\Large\sffamily\bfseries\color{DeepNavy}}{\thesection}{0.8em}{}
\titleformat{\subsection}{\large\sffamily\bfseries\color{Teal}}{\thesubsection}{0.8em}{}
\titleformat{\subsubsection}{\normalsize\sffamily\bfseries\color{DarkGray}}{\thesubsubsection}{0.8em}{}

\pagestyle{fancy}
\fancyhf{}
\fancyhead[L]{\small\sffamily カオス同期・情報圧縮 数値実験テキスト}
\fancyhead[R]{\small\sffamily\leftmark}
\fancyfoot[C]{\small\thepage}
\renewcommand{\headrulewidth}{0.3pt}
\renewcommand{\footrulewidth}{0pt}

% ============================================================
% 記号・環境
% ============================================================
\newcommand{\R}{\mathbb{R}}
\newcommand{\C}{\mathbb{C}}
\newcommand{\E}{\mathbb{E}}
\newcommand{\Var}{\operatorname{Var}}
\newcommand{\median}{\operatorname{median}}
\newcommand{\MAD}{\operatorname{MAD}}
\newcommand{\IQR}{\operatorname{IQR}}
\newcommand{\diag}{\operatorname{diag}}
\newcommand{\rank}{\operatorname{rank}}
\newcommand{\tr}{\operatorname{tr}}
\newcommand{\Proj}{\operatorname{Proj}}
\newcommand{\NMSE}{\ensuremath{\operatorname{NMSE}}}
\newcommand{\PSNR}{\ensuremath{\operatorname{PSNR}}}
\newcommand{\SSIM}{\ensuremath{\operatorname{SSIM}}}
\newcommand{\effrank}{\ensuremath{\operatorname{erank}}}
\newcommand{\dd}{\,\mathrm{d}}
\newcommand{\e}{\mathrm{e}}
\newcommand{\ii}{\mathrm{i}}
\newcommand{\ones}{\bm{1}}
\newcommand{\vect}[1]{\bm{#1}}
\newcommand{\mat}[1]{\bm{#1}}
\newcommand{\code}[1]{\texttt{#1}}
\newcommand{\EID}[1]{\textsf{#1}}
\newcommand{\Gate}{\textsf{Gate}}

\newtcolorbox{purposebox}[1][]{
  enhanced,breakable,colback=LightBlue,colframe=DeepNavy,boxrule=0.8pt,
  title={この章の目的},fonttitle=\sffamily\bfseries,#1}
\newtcolorbox{theorybox}[1][]{
  enhanced,breakable,colback=LightTeal,colframe=Teal,boxrule=0.7pt,
  title={理論の要点},fonttitle=\sffamily\bfseries,#1}
\newtcolorbox{cautionbox}[1][]{
  enhanced,breakable,colback=LightRed,colframe=Brick,boxrule=0.8pt,
  title={重要な注意},fonttitle=\sffamily\bfseries,#1}
\newtcolorbox{taskbox}[1][]{
  enhanced,breakable,colback=LightOrange,colframe=Orange,boxrule=0.8pt,
  title={実験課題},fonttitle=\sffamily\bfseries,#1}
\newtcolorbox{gatebox}[1][]{
  enhanced,breakable,colback=LightGray,colframe=Forest,boxrule=0.8pt,
  title={事前に固定する合否条件},fonttitle=\sffamily\bfseries,#1}
\newtcolorbox{auditbox}[1][]{
  enhanced,breakable,colback=white,colframe=Purple,boxrule=0.8pt,
  title={主張監査},fonttitle=\sffamily\bfseries,#1}
\newtcolorbox{codebox}[1][]{
  enhanced,breakable,colback=LightGray,colframe=MidGray,boxrule=0.5pt,
  title={実装契約},fonttitle=\sffamily\bfseries,#1}
\newtcolorbox{reportbox}[1][]{
  enhanced,breakable,colback=white,colframe=DeepNavy,boxrule=0.7pt,
  title={レポートに必ず書くこと},fonttitle=\sffamily\bfseries,#1}

\newcommand{\Known}{\textcolor{Forest}{\sffamily\bfseries [Known]}}
\newcommand{\SourceClaim}{\textcolor{Brick}{\sffamily\bfseries [Source claim]}}
\newcommand{\Derived}{\textcolor{Teal}{\sffamily\bfseries [Derived]}}
\newcommand{\Numerical}{\textcolor{Purple}{\sffamily\bfseries [Numerical]}}
\newcommand{\Hypothesis}{\textcolor{Orange}{\sffamily\bfseries [Hypothesis]}}

\lstdefinestyle{pythonstyle}{
  language=Python,
  basicstyle=\small\ttfamily,
  keywordstyle=\bfseries\color{DeepNavy},
  commentstyle=\itshape\color{Forest},
  stringstyle=\color{Brick},
  showstringspaces=false,
  columns=fullflexible,
  keepspaces=true,
  frame=single,
  framerule=0.3pt,
  rulecolor=\color{MidGray},
  backgroundcolor=\color{LightGray},
  breaklines=true,
  tabsize=4,
  numbers=left,
  numberstyle=\tiny\color{MidGray},
  xleftmargin=2em,
  framexleftmargin=1.5em
}

\crefname{equation}{式}{式}
\crefname{figure}{図}{図}
\crefname{table}{表}{表}
\crefname{chapter}{第}{第}
\crefname{section}{節}{節}

\begin{document}

% ============================================================
% 表紙
% ============================================================
\begin{titlepage}
  \thispagestyle{empty}
  \begin{tikzpicture}[remember picture,overlay]
    \fill[DeepNavy] (current page.north west) rectangle ([yshift=-52mm]current page.north east);
    \fill[Teal] ([yshift=-52mm]current page.north west) rectangle ([yshift=-57mm]current page.north east);
    \draw[DeepNavy,line width=1.2pt] ([xshift=18mm,yshift=18mm]current page.south west)
      rectangle ([xshift=-18mm,yshift=-18mm]current page.north east);
  \end{tikzpicture}

  \vspace*{17mm}
  {\color{white}\sffamily\bfseries\fontsize{26}{34}\selectfont
    カオス同期とTM読み出しによる\\[2mm]
    情報圧縮}

  \vspace{9mm}
  {\color{white}\sffamily\fontsize{17}{22}\selectfont
    数値実験テキスト - E0からE5までを自分で実装するために}

  \vspace{42mm}
  \begin{center}
    {\sffamily\bfseries\Large 一般化Boole写像・Cauchy安定分布・Cayley変換}\\[2mm]
    {\sffamily\bfseries\Large Takenaka--Malmquist基底・同期転移・復号可能性}\\[17mm]

    \begin{tcolorbox}[width=0.84\textwidth,colback=LightBlue,colframe=DeepNavy,boxrule=0.8pt]
    \centering
    本書の中心課題は、\textbf{同期による自由度縮約}と\textbf{入力情報を保持した圧縮}を分離して測り、\par
    「どの条件なら両立し、どの条件では情報が潰れるか」を反証可能な実験へ落とすことである。
    \end{tcolorbox}

    \vfill
    {\large 長松蒼空}\\[2mm]
    {\small 研究用実験テキスト v1.0}\\[1mm]
    {\small 2026年9月4日}
  \end{center}
\end{titlepage}

\frontmatter

\chapter*{本書の位置づけ}
\addcontentsline{toc}{chapter}{本書の位置づけ}

本書は、添付された「2024年度数理工学実験テキスト」の構成原則を参考に、カオス同期と情報圧縮に関する研究を、自分で数式を確認し、自分で実装し、自分で結果を解釈できるように再構成した実験書である。各実験は、\textbf{目的、原理、実験方法、保存物、評価指標、合否条件、考察課題}の順に記述する。科学的文書は、実験を知らない第三者が再現できる情報を含み、式・図・表を番号で参照し、結果と主観的感想を分ける、という方針を採用する\cite{mathlab2024}。

研究の最上位目的は次である。

\begin{quote}
\textbf{神経系に見られる動的同期に着想を得て、カオス同期を情報圧縮へ応用できるかを検討する。}
\end{quote}

ただし、「同期した」ことだけを圧縮成功と呼ばない。完全同期により全入力が同じ軌道へ潰れた場合、幾何学的な自由度は減っていても、復号器に必要な情報は失われている。本書では、同期、局所安定性、入力識別、復元誤差、潜在次元を別々に測る。

\begin{cautionbox}
本書は、添付原稿に書かれた強い理論主張を無条件に事実として扱わない。特に、無限次元ランダム結合の厳密縮約、Perron--Frobenius固有関数、KMS状態、質量ギャップに関する主張は、\SourceClaim として区別する。卒業研究で確定事項として書けるのは、標準理論、独立導出、または自分の再現計算で確認した範囲である。
\end{cautionbox}

\section*{主張ラベル}
本書では各主張の確度を次のように表す。

\begin{longtable}{p{32mm}p{112mm}}
\toprule
ラベル & 意味 \\
\midrule
\Known & 標準的な数学、または信頼できる原典で確認できる事項。\\
\SourceClaim & 添付資料に記載されているが、本研究では独立監査が必要な事項。\\
\Derived & 本書で仮定を明示して導出した事項。\\
\Numerical & 既存runまたは本実装で数値的に確認する事項。\\
\Hypothesis & 今後の実験で支持・棄却する研究仮説。\\
\bottomrule
\end{longtable}

\section*{使い方}
\begin{enumerate}
  \item 第\ref{chap:practice}章で、レポートと再現性の契約を先に固定する。
  \item 第\ref{chap:foundation}章で、Cauchy分布、一般化Boole写像、Cayley変換、TM基底を紙と鉛筆で確認する。
  \item \EID{E0A}から順に実装し、各章末の\Gate を通過するまで次段階へ進まない。
  \item 数値結果が理論に合わなくても、clipやfallbackで隠さず、失敗条件を保存する。
  \item 各実験終了後に「結果」と「考察」を同じrunフォルダへ保存する。
\end{enumerate}

\tableofcontents

\mainmatter

% ============================================================
\chapter{実験の進め方と再現性契約}
\label{chap:practice}
% ============================================================

\begin{purposebox}
研究コードを書く前に、何を仮説とし、何を固定し、何を出力し、どの結果なら次へ進むかを決める。本章終了時には、一つのrunを第三者が再実行できる最小契約を説明できることを目標とする。
\end{purposebox}

\section{実験レポートの構成}

本研究のレポートは、数理工学実験テキストの推奨構成に従い、少なくとも次の七項目を含める\cite{mathlab2024}。

\begin{enumerate}
  \item \textbf{はじめに}: 何を確かめたいか、なぜその問いが必要かを書く。
  \item \textbf{原理・理論}: 使用する写像、分布、同期条件、評価量を自分の言葉で導出する。
  \item \textbf{実験方法}: データ生成、パラメータ、seed、分割、前処理、アルゴリズムを再現可能に書く。
  \item \textbf{結果}: 数値、図、表を提示する。採否判断はまだ混ぜない。
  \item \textbf{考察}: 理論との一致・不一致、代替説明、数値誤差、限界を議論する。
  \item \textbf{結論}: 何が支持され、何が未支持かを短くまとめる。
  \item \textbf{参考文献}: どの主張が誰の結果に依存するかを明示する。
\end{enumerate}

\begin{reportbox}
「同期誤差が小さくなった」という結果だけでは不十分である。少なくとも、同期誤差、入力復号性能、潜在次元、seed間分散を同じ条件表に載せる。理論値との比較では、式の仮定と実装したモデルが同じかを先に確認する。
\end{reportbox}

\section{一つのrunの標準構造}

実験IDは研究上の依存関係を表す。旧E/T系列は、次の正規系列に統一する。

\begin{table}[htbp]
\centering
\caption{正規実験系列。E0からE5へ進むほど応用側へ移る。}
\label{tab:canonical_stages}
\begin{tabularx}{\textwidth}{>{\bfseries}p{14mm}p{43mm}X}
\toprule
段階 & 研究上の役割 & この段階で分離する失敗原因 \\
\midrule
E0 & 単一写像の数学・数値整合性 & 写像、分布、特異点処理、理論式の誤り \\
E1 & TM時間特徴の読み出し・識別可能性 & 特徴量の不足、観測写像による情報損失 \\
E2 & 雑音・欠測・有限長への頑健性 & 前処理と基底の脆弱性、有限標本誤差 \\
E3 & 同期転移・basin・入力保持 & 局所安定性と大域挙動、同期と情報消失 \\
E4 & 合成時系列の圧縮・復元 & 潜在座標と復号器、rate--distortion \\
E5 & 画像の圧縮・復元 & 実データへの転移、baselineとの公平比較 \\
\bottomrule
\end{tabularx}
\end{table}

正規IDは、例えば
\begin{center}
\code{E3B-TANGENT-TWO-NODE-BASIN},\quad
\code{E3C-BOOLE-N8-INPUT-RETENTION}
\end{center}
のように、段階、写像族、対象、目的が名前から分かるようにする。実行フォルダは
\begin{center}
\code{YYYYMMDD\_<正規ID>\_<slug>}
\end{center}
とする。旧runは参照整合性やハッシュを壊すため改名せず、レジストリで正規IDへ対応付ける。

\begin{figure}[htbp]
\centering
\begin{tikzpicture}[
  node distance=8mm and 3.5mm,
  stage/.style={rounded corners,draw=DeepNavy,fill=LightBlue,minimum width=22mm,minimum height=12mm,align=center,font=\sffamily\bfseries\small},
  arrow/.style={-{Latex[length=2.6mm]},line width=0.8pt,DeepNavy}
]
\node[stage] (e0) {E0\\基礎力学};
\node[stage,right=of e0] (e1) {E1\\読み出し};
\node[stage,right=of e1] (e2) {E2\\頑健性};
\node[stage,right=of e2] (e3) {E3\\同期と情報};
\node[stage,right=of e3] (e4) {E4\\信号復元};
\node[stage,right=of e4] (e5) {E5\\画像pilot};
\draw[arrow] (e0)--(e1);
\draw[arrow] (e1)--(e2);
\draw[arrow] (e2)--(e3);
\draw[arrow] (e3)--(e4);
\draw[arrow] (e4)--(e5);
\node[below=7mm of e2,align=center,font=\small] {各矢印は「前段の\Gate を通過してから進む」を意味する};
\end{tikzpicture}
\caption{実験の依存関係。いきなり画像へ進まず、失敗原因を一段ずつ切り分ける。}
\label{fig:experiment_dependency}
\end{figure}

各runには最低限、次を保存する。

\begin{longtable}{p{45mm}p{98mm}}
\toprule
保存物 & 内容 \\
\midrule
\code{config.yaml} & 実験前に固定した全パラメータ、seed、データ分割、\Gate。\\
\code{environment.json} & Python、NumPy、SciPy、scikit-learn、PyTorch等のバージョンと計算機情報。\\
\code{metrics.csv} & 条件別・seed別の未集約指標。平均だけを保存しない。\\
\code{predictions.npz} & 全test予測、必要なら潜在表現とターゲット。\\
\code{summary.json} & 集約値、信頼区間、採否、失敗理由。\\
\code{figures/} & 図と、その元データ。画像だけを数値正本にしない。\\
\code{validation.json} & 単体テスト、理論整合性テスト、有限値チェック。\\
\code{sha256.txt} & 主要artifactのハッシュ。\\
\path{experiment_note.md} & 目的、方法、結果、考察、次の一手。\\
\bottomrule
\end{longtable}

\section{データ漏洩を防ぐ順序}

前処理、PCA、ridge係数、閾値、結合強度の選択は、すべてtrainまたはvalidationだけで行う。testは最後に一度だけ評価する。標準手順は次である。

\begin{enumerate}
  \item seedとサンプル生成単位を決め、group単位でtrain/validation/testへ分ける。
  \item trainだけでrobust scaler、PCA、decoderをfitする。
  \item validationだけで力学パラメータや潜在次元候補を選ぶ。
  \item 選択を固定し、testで一度だけ最終性能を測る。
  \item testを見て選択を変えた場合は、そのtestをvalidationへ降格し、新しいtestを作る。
\end{enumerate}

\begin{cautionbox}
同じ軌道から切り出した窓をランダムに分割すると、初期条件や長時間相関がtrainとtestに共有される。軌道seed、パラメータ値、元画像など、依存を生む最上位単位で分割する。
\end{cautionbox}

\section{実装規約}

本プロジェクトでは、コードの責務を次のように分離する。

\begin{itemize}
  \item \code{maps.py}: 局所写像と解析式。
  \item \code{coupling.py}: 結合行列、同期多様体、ネットワーク更新。
  \item \code{features.py}: Cayley/TM、Fourier、graph mode、時間相関。
  \item \code{metrics.py}: 同期誤差、Lyapunov推定、復元指標、effective rank。
  \item \code{datasets.py}: 合成信号、二源観測、digits系列化。
  \item \code{experiments/}: 各EIDの設定とrunner。
\end{itemize}

関数・クラスのdocstringには\textbf{How}、テスト名・docstringには\textbf{What}、コミットメッセージには\textbf{Why}、コードコメントには\textbf{Why not}を残す。例えば、特異点を単純clipしない理由は、clipにより元の写像とは異なる力学系になるためである。

\begin{codebox}
大規模runnerの前に、次の小テストを通す。
\begin{enumerate}
  \item 理論固定点を代入した残差が十分小さい。
  \item 同期状態を入力すると同期状態のまま更新される。
  \item TMモードの絶対値が浮動小数誤差の範囲で1である。
  \item train統計がtestへ混入していない。
  \item 非有限値、発散、特異点近傍通過率を数える。
\end{enumerate}
\end{codebox}

\section{完了の定義}

「completed」は仮説支持を意味しない。設定、全予測、未集約指標、環境、検証、採否理由が揃ったことを意味する。仮説が棄却されても、再現可能な負の結果なら完了である。一方、良い図が一枚できてもseed別結果や設定が失われていれば未完了である。

\begin{taskbox}
\textbf{課題 1.1}\quad 付属テンプレートを用い、\EID{E0A-BOOLE-LOCAL-VALIDATION}の\code{config.yaml}を作成せよ。まだコードは書かず、目的変数、操作変数、固定変数、出力、\Gate、失敗条件を文章化すること。

\medskip
\textbf{課題 1.2}\quad 自分の環境で\code{environment.json}を生成し、同じコードがGoogle ColabとKaggle Notebookでどの程度同じ結果を返すべきか、乱数・BLAS・GPU演算の観点から述べよ。
\end{taskbox}

% ============================================================
\chapter{数学的準備: Cauchy分布・Boole写像・TM基底}
\label{chap:foundation}
% ============================================================

\begin{purposebox}
重い裾を持つCauchy分布、一般化Boole写像の不変尺度、Cayley変換、TM power basisを理解する。E0以降で使う数式を、実装可能な形まで導く。
\end{purposebox}

\section{Cauchy分布}

位置\(\mu\in\R\)、尺度\(\gamma>0\)のCauchy分布を
\begin{equation}
  p_{\mu,\gamma}(x)
  = \frac{1}{\pi}\frac{\gamma}{(x-\mu)^2+\gamma^2}
  \label{eq:cauchy_density}
\end{equation}
で定義する。その特性関数は
\begin{equation}
  \varphi_{\mu,\gamma}(t)
  = \E[\e^{\ii tX}]
  = \exp\!\left(\ii\mu t-\gamma|t|\right)
  \label{eq:cauchy_cf}
\end{equation}
である。

\Known \cref{eq:cauchy_cf}から、独立な\(X_j\sim\mathrm{Cauchy}(\mu_j,\gamma_j)\)と実係数\(a_j\)に対し、
\begin{equation}
  \sum_{j=1}^{m}a_jX_j
  \sim
  \mathrm{Cauchy}\!\left(
    \sum_{j=1}^{m}a_j\mu_j,
    \sum_{j=1}^{m}|a_j|\gamma_j
  \right)
  \label{eq:cauchy_linear_stability}
\end{equation}
が成り立つ。これがCauchy分布の安定性である。

Cauchy分布では通常の平均と分散が存在しない。そのため、実験では標本平均・標本分散を尺度推定の中心に置かない。中心にはmedian、尺度には
\begin{equation}
  \widehat\gamma_{\mathrm{MAD}}
  = \median_t |x_t-\median_s x_s|
  \label{eq:cauchy_mad}
\end{equation}
またはhalf-IQRを用いる。理想的なCauchy分布では、MADとhalf-IQRはいずれも\(\gamma\)に等しい。

\begin{cautionbox}
標本平均が有限の値を返しても、母平均が存在することを意味しない。軌道長を増やしたときに標本平均が安定しないこと自体が、実装の異常ではなくCauchy裾の性質である可能性がある。
\end{cautionbox}

\section{一般化Boole写像}

本研究の基本局所写像を
\begin{equation}
  F_{\alpha}(x)
  = \alpha\left(x-\frac{1}{x}\right),
  \qquad 0<\alpha<1
  \label{eq:boole_map}
\end{equation}
とする。\(x=0\)は特異点である。

\subsection{Cauchy尺度の更新}

\SourceClaim 添付の無限次元同期原稿では、\(X\sim\mathrm{Cauchy}(0,\gamma)\)に\cref{eq:boole_map}を作用させると、像の分布もCauchy型に保たれ、尺度が
\begin{equation}
  G_{\alpha}(\gamma)
  = \alpha\left(\gamma+\frac{1}{\gamma}\right)
  \label{eq:boole_scale_map}
\end{equation}
へ更新されるとする\cite{infinite_sync}。この写像の正固定点は、
\begin{align}
  \gamma_* &= \alpha\left(\gamma_*+\frac{1}{\gamma_*}\right),\\
  (1-\alpha)\gamma_*^2 &= \alpha,
\end{align}
より
\begin{equation}
  \gamma_*
  = \sqrt{\frac{\alpha}{1-\alpha}}
  \label{eq:boole_gamma_star}
\end{equation}
である。尺度写像の微分は
\begin{equation}
  G'_{\alpha}(\gamma)
  = \alpha-\frac{\alpha}{\gamma^2}
\end{equation}
であり、固定点では
\begin{equation}
  G'_{\alpha}(\gamma_*)=2\alpha-1.
  \label{eq:scale_stability}
\end{equation}
したがって\(0<\alpha<1\)で\(|2\alpha-1|<1\)となり、尺度固定点は局所的に吸引的である。

\subsection{Lyapunov指数}

写像の微分は
\begin{equation}
  F'_{\alpha}(x)
  = \alpha\left(1+\frac{1}{x^2}\right)
  \label{eq:boole_derivative}
\end{equation}
である。不変Cauchy測度を仮定すると、Lyapunov指数は
\begin{equation}
  \lambda_F(\alpha)
  = \int_{-\infty}^{\infty}
    p_{0,\gamma_*}(x)
    \log\left|F'_{\alpha}(x)\right|\dd x
  \label{eq:boole_lyapunov_integral}
\end{equation}
で定義される。添付原稿の留数積分を用いると
\begin{align}
  \lambda_F(\alpha)
  &= \log\alpha
     +2\log\left(1+\frac{1}{\gamma_*}\right)\\
  &=2\log\left(\sqrt\alpha+\sqrt{1-\alpha}\right)
  \label{eq:boole_lyapunov_closed}
\end{align}
となる\cite{infinite_sync}。特に\(\alpha=1/2\)では
\begin{equation}
  \gamma_*=1,\qquad \lambda_F=\log 2.
  \label{eq:maxchaos_boole}
\end{equation}

\begin{auditbox}
\cref{eq:boole_scale_map,eq:boole_lyapunov_closed}は、E0Aで数値再現する対象である。コードが式をそのまま返すだけでは検証にならない。独立な軌道生成からMADと時間平均Lyapunov指数を推定し、理論値へ近づくかを調べる。
\end{auditbox}

\section{Cayley変換とTM power basis}

複素極\(z_0=\mu+\ii\gamma\)に対し、Cayley変換
\begin{equation}
  q_{\mu,\gamma}(x)
  = \frac{(x-\mu)-\ii\gamma}{(x-\mu)+\ii\gamma}
  \label{eq:cayley}
\end{equation}
を定義する。実数\(x\)では分子と分母が共役なので
\begin{equation}
  |q_{\mu,\gamma}(x)|=1.
  \label{eq:cayley_modulus}
\end{equation}

角度変数\(\theta\in(0,2\pi)\)を
\begin{equation}
  x=\mu+\gamma\cot\frac{\theta}{2}
  \label{eq:cauchy_angle}
\end{equation}
で導入すると、向きを除いて
\begin{equation}
  p_{\mu,\gamma}(x)|\dd x|
  =\frac{\dd\theta}{2\pi},
  \qquad
  q_{\mu,\gamma}(x)=\e^{-\ii\theta}.
  \label{eq:cauchy_uniform_circle}
\end{equation}
したがって、Cauchy測度付き実直線は一様測度付き単位円へ移る。

TM power basisを
\begin{equation}
  M_k^{(\mu,\gamma)}(x)
  =q_{\mu,\gamma}(x)^k,
  \qquad k\in\mathbb Z
  \label{eq:tm_power_basis}
\end{equation}
と定義する。\cref{eq:cauchy_uniform_circle}より
\begin{equation}
  \left\langle M_k,M_{\ell}\right\rangle
  =\int_{\R}\overline{M_k(x)}M_{\ell}(x)
    p_{\mu,\gamma}(x)\dd x
  =\delta_{k\ell}.
  \label{eq:tm_orthonormal}
\end{equation}

\begin{figure}[htbp]
\centering
\begin{tikzpicture}[
  box/.style={rounded corners,draw=DeepNavy,fill=LightBlue,text width=30mm,minimum height=14mm,align=center,font=\sffamily\bfseries\small},
  arrow/.style={-{Latex[length=2.6mm]},line width=0.9pt,Teal}
]
\node[box] (x) {重い裾の軌道\\\(x_t\in\R\)};
\node[box,right=6mm of x] (q) {Cayley変換\\\(q_t\in\mathbb S^1\)};
\node[box,right=6mm of q] (m) {TMモード\\\(M_k(t)=q_t^k\)};
\node[box,right=6mm of m] (f) {時間相関・graph射影\\有限次元特徴};
\draw[arrow] (x)--(q);
\draw[arrow] (q)--(m);
\draw[arrow] (m)--(f);
\node[below=8mm of q,align=center,font=\small] {巨大な\(|x_t|\)を直接使わず、\(|q_t|=1\)の位相へ写す};
\end{tikzpicture}
\caption{Cauchy軌道からTM時間特徴を作る流れ。}
\label{fig:cauchy_tm_pipeline}
\end{figure}

\section{Koopman作用が厳密に単純化する基準系}

\SourceClaim 添付のK-fold原稿で、本研究に直接必要なのはKMS状態や質量ギャップではなく、K-fold cotangent写像
\begin{equation}
  H_K(x)=\cot\left(K\operatorname{arccot}x\right),
  \qquad K=2,3,\dots
  \label{eq:kfold_cotangent}
\end{equation}
とTM基底の対応である\cite{mass_gap}。標準Cauchy分布\((\mu,\gamma)=(0,1)\)に対し、Koopman作用素\(U_H f=f\circ H_K\)は
\begin{equation}
  U_H M_k=M_{Kk}
  \label{eq:koopman_shift}
\end{equation}
と基底添字を移す。非線形写像が整数添字のシフトとして見えるため、実装の正対照に適している。

\begin{cautionbox}
\cref{eq:koopman_shift}が厳密に成り立つのは、\cref{eq:kfold_cotangent}のようなK-fold基準系である。任意の一般化Boole写像\(F_\alpha\)や、生の\(\tan(\beta x)\)へ同じ式をそのまま流用しない。また、K-fold tangent写像\(\tan(K\arctan x)\)と、生のtangent写像\(\tan(\beta x)\)は別の力学系である。
\end{cautionbox}

\section{時間相関としての読み出し}

単一時刻のTM係数だけでなく、lag\(\ell\)の複素自己相関
\begin{equation}
  C_{k}(\ell)
  =\frac{1}{T-\ell}
   \sum_{t=0}^{T-\ell-1}
   \overline{M_k(x_t)}M_k(x_{t+\ell})
  \label{eq:tm_autocorrelation}
\end{equation}
を用いる。\(\Re C_k(\ell)\)と\(\Im C_k(\ell)\)を特徴とすれば、周辺分布をrobust標準化した後にも残る時間順序を読むことができる。

\Hypothesis 本研究の中心的な読み出し仮説は、「TM基底が情報を作る」ではなく、\textbf{軌道に残っている入力依存の時間構造を、重い裾に頑健な有界座標で読む}ことである。観測写像が情報を捨てた場合、TM変換はそれを復元できない。

\begin{taskbox}
\textbf{課題 2.1}\quad \cref{eq:cauchy_cf}を用いて\cref{eq:cauchy_linear_stability}を導け。

\medskip
\textbf{課題 2.2}\quad \cref{eq:cauchy_angle}を微分し、\cref{eq:cauchy_uniform_circle}を確認せよ。また\cref{eq:tm_orthonormal}をFourier直交性へ帰着せよ。

\medskip
\textbf{課題 2.3}\quad \(\alpha\in\{0.25,0.5,0.75\}\)について、\(\gamma_*\)と\(\lambda_F\)を計算せよ。\(\lambda_F\)が\(\alpha=1/2\)で最大になることを微分または対称性から説明せよ。
\end{taskbox}

% ============================================================
\chapter{E0: 単一写像と基底の数値検証}
\label{chap:e0}
% ============================================================

\begin{purposebox}
結合系へ進む前に、局所写像、特異点処理、不変Cauchy尺度、Lyapunov指数、TM基底の実装が理論と整合することを確認する。E0を通過しないコードで同期や圧縮を議論してはならない。
\end{purposebox}

\section{E0A-BOOLE-LOCAL-VALIDATION}

\subsection{研究質問}

\begin{quote}
一般化Boole写像\(F_\alpha\)を浮動小数点で反復したとき、burn-in後の単一軌道は理論Cauchy尺度とLyapunov指数を再現するか。
\end{quote}

操作変数は\(\alpha\)、軌道長、初期分布、浮動小数精度である。主要出力は、尺度推定誤差、Lyapunov誤差、非有限値率、特異点近傍通過率である。

\subsection{推奨条件}

\begin{table}[htbp]
\centering
\caption{E0Aの標準条件。最初は小さく動かし、合格後に本条件へ増やす。}
\label{tab:e0a_config}
\begin{tabularx}{\textwidth}{p{37mm}X}
\toprule
項目 & 推奨値 \\
\midrule
\(\alpha\) & \(0.25,0.40,0.50,0.60,0.75\)。\(0.5\)を基準条件とする。\\
seed & 最低5。初期値をCauchy、Gaussian、一様分布から生成する。\\
burn-in & まず\(2\times10^4\)。尺度収束が遅い条件では増やす。\\
標本長 & pilot \(10^4\)、本実験 \(10^5\)以上。\\
精度 & float64を主、float32を監査。可能ならlong doubleも比較。\\
特異点近傍 & \(|x|<10^{-12},10^{-10},10^{-8}\)の通過回数を保存。\\
\bottomrule
\end{tabularx}
\end{table}

\subsection{実験手順}

\begin{enumerate}
  \item \(x_0\)を生成し、\cref{eq:boole_map}をburn-in回反復する。
  \item burn-in後の軌道\(x_0,\dots,x_{T-1}\)を保存する。
  \item median、MAD、half-IQRを計算し、\cref{eq:boole_gamma_star}と比較する。
  \item \cref{eq:boole_derivative}から
  \begin{equation}
    \widehat\lambda_F
    =\frac1T\sum_{t=0}^{T-1}
      \log\left|\alpha\left(1+x_t^{-2}\right)\right|
    \label{eq:empirical_boole_lyap}
  \end{equation}
  を計算する。
  \item 理論密度とのヒストグラムだけでなく、Cauchy QQプロットを作る。
  \item 軌道長を増やしたときの誤差曲線を描く。
\end{enumerate}

\begin{codebox}
\textbf{How}: 更新関数は状態配列を受け取り、非有限値マスクと特異点近傍マスクを同時に返す。\par
\textbf{Why not}: \(x=0\)付近を固定値へclipして計算を続けると、元のBoole写像とは異なる写像になる。発散や非有限値は失敗として記録し、clip版を使う場合は別EIDにする。
\end{codebox}

\subsection{図と表}

必須図は次の四つである。

\begin{enumerate}
  \item 理論Cauchy密度と経験密度。
  \item Cauchy QQプロット。
  \item 理論\(\gamma_*\)対推定MADの散布図。破線は完全一致。
  \item 理論\(\lambda_F\)対\(\widehat\lambda_F\)の散布図。破線は完全一致。
\end{enumerate}

表には、\(\alpha\)、seed、\(\gamma_*\)、\(\widehat\gamma\)、絶対誤差、\(\lambda_F\)、\(\widehat\lambda_F\)、非有限値率、実行時間を載せる。

\begin{gatebox}
\begin{enumerate}
  \item 各\(\alpha\)でMAD相対誤差のseed中央値が5\%以内。
  \item Lyapunov絶対誤差のseed中央値が\(2\times10^{-2}\)以内。
  \item 軌道長を増やすと誤差が全体として減少する。
  \item float32とfloat64で結論が反転しない。反転する場合は原因を特定する。
  \item 数値安定化が元写像を変更していない。
\end{enumerate}
一つでも満たさなければ、結合系へ進まず、特異点処理、burn-in、引数型、乱数生成を監査する。
\end{gatebox}

\section{E0B-KFOLD-TM-SHIFT-CHECK}

\subsection{目的}

K-fold cotangent写像を基準系とし、Cauchy測度、Cayley変換、TM直交性、Koopman shiftを数値的に確認する。これは一般化Boole写像のすべての\(\alpha\)に対する証明ではなく、TM実装の正対照である。

\subsection{手順}

\begin{enumerate}
  \item \(\theta_t\sim\mathrm{Uniform}(0,\pi)\)から\(x_t=\cot\theta_t\)を生成する。
  \item \(K\in\{2,3,5\}\)、\(k\in\{1,2,3,4\}\)で\(H_K(x)\)と\(M_k(x)\)を実装する。
  \item 各標本について
  \begin{equation}
    r_{K,k}(x)
    =\left|M_k(H_K(x))-M_{Kk}(x)\right|
    \label{eq:tm_shift_residual}
  \end{equation}
  を計算する。
  \item Monte Carlo内積
  \begin{equation}
    \widehat G_{k\ell}
    =\frac1T\sum_{t=1}^{T}\overline{M_k(x_t)}M_\ell(x_t)
    \label{eq:tm_gram_empirical}
  \end{equation}
  を作り、単位行列との差を測る。
\end{enumerate}

\begin{cautionbox}
\(\operatorname{arccot}\)の枝と\(\cot(K\operatorname{arccot}x)\)の実装はライブラリにより異なる。数値的には\(\theta=\operatorname{atan2}(1,x)\in(0,\pi)\)を用い、\(\cot(K\theta)=\cos(K\theta)/\sin(K\theta)\)として枝を固定する。
\end{cautionbox}

\begin{gatebox}
\begin{enumerate}
  \item \(\max r_{K,k}\)がfloat64の許容誤差内である。ただし極近傍は別集計する。
  \item \(\|\widehat G-I\|_{\mathrm{F}}\)が標本長とともに減少する。
  \item 乱数seedを変えても直交性の結論が変わらない。
\end{enumerate}
\end{gatebox}

\section{E0の考察課題}

\begin{taskbox}
\textbf{課題 3.1}\quad 標本平均、標本標準偏差、MADを軌道長に対して描け。Cauchy分布において、なぜMADだけが安定するか説明せよ。

\medskip
\textbf{課題 3.2}\quad \(\alpha\to0\)または\(\alpha\to1\)で\(\gamma_*\)と\(\lambda_F\)がどう振る舞うか調べ、数値実験が難しくなる理由を述べよ。

\medskip
\textbf{課題 3.3}\quad 特異点近傍をclipした実装と、失敗を記録して停止する実装を比較せよ。密度、Lyapunov指数、再現性にどのような差が生じるかを示せ。
\end{taskbox}

% ============================================================
\chapter{E1: TM時間特徴と識別可能性}
\label{chap:e1}
% ============================================================

\begin{purposebox}
TM特徴が、有限軌道に残るパラメータ情報を読めるかを検証する。同時に、観測写像が情報を失った場合には、どの特徴量でも復元できないという識別可能性境界を正対照・負対照で確認する。
\end{purposebox}

\section{E1A-BOOLE-SINGLE-SOURCE-READOUT}

\subsection{研究質問}

周辺Cauchy尺度をrobust標準化して除いた後にも、有限時間の軌道順序には\(\alpha\)依存情報が残るか。その情報を、同じ特徴次元のTM時間相関とFourier特徴のどちらが安定に読み出せるか。

\subsection{データ生成}

trainとtestで\(\alpha\)値を分け、補間能力を測る。推奨例は
\begin{align}
  \mathcal A_{\mathrm{train}}
  &=\{0.34,0.42,0.50,0.58,0.66\},\\
  \mathcal A_{\mathrm{test}}
  &=\{0.38,0.46,0.54,0.62\}.
  \label{eq:e1_alpha_split}
\end{align}
各\(\alpha\)について独立seedで複数軌道を生成する。軌道内で
\begin{equation}
  \widetilde x_t
  =\frac{x_t-\median_s x_s}{\IQR(x)/2}
  \label{eq:robust_orbit_standardization}
\end{equation}
と標準化し、周辺尺度だけで\(\alpha\)を当てる近道を抑える。

\subsection{TM特徴}

標準Cauchy極を用いて
\begin{equation}
  q_t=\frac{\widetilde x_t-\ii}{\widetilde x_t+\ii},
  \qquad M_k(t)=q_t^k
  \label{eq:e1_tm_modes}
\end{equation}
とし、\(k=1,\dots,4\)、\(\ell\in\{1,2\}\)について\cref{eq:tm_autocorrelation}を計算する。実部・虚部を連結すると
\begin{equation}
  4\text{ modes}\times2\text{ lags}\times2
  =16
  \label{eq:e1_tm_dim}
\end{equation}
次元となる。

比較するFourier特徴も16次元へそろえる。例えば、窓関数を掛けたrFFTのパワーを16帯域に分け、各帯域の\(\log(\epsilon+\text{energy})\)を用いる。

\subsection{負対照}

時間順序を本当に読んでいるかを確認するため、次を入れる。

\begin{enumerate}
  \item \textbf{shuffle対照}: 各軌道の時刻順序をランダムに並べ替えて同じ特徴を作る。
  \item \textbf{lagなし対照}: \(T^{-1}\sum_tM_k(t)\)のみを使う。
  \item \textbf{quantile対照}: median、IQR、複数分位点だけを使う。
\end{enumerate}

\subsection{推定器と評価}

同じ正則化候補を持つridge回帰で\(\alpha\)を推定する。正則化係数はtrain内CVまたはvalidationで選ぶ。主要指標はRMSE、MAE、\(R^2\)、seedごとの誤差差である。

TMとFourierのpaired誤差差を
\begin{equation}
  \Delta_s
  =\mathrm{RMSE}_{\mathrm{TM},s}
   -\mathrm{RMSE}_{\mathrm{Fourier},s}
  \label{eq:paired_rmse_difference}
\end{equation}
とし、seed単位bootstrapで95\%信頼区間を作る。

\begin{gatebox}
\begin{enumerate}
  \item TMまたはFourierの少なくとも一方が、shuffle対照より有意に良い。
  \item trainにない\(\alpha\)で性能を報告する。
  \item 特徴次元、decoder、正則化探索予算をそろえる。
  \item TM優位を主張する場合、\(\Delta\)の95\%上限が0未満である。
  \item lagなしTMが同等なら、「時間相関が必要」という仮説は未支持とする。
\end{enumerate}
\end{gatebox}

\begin{reportbox}
既存runの参考値として、16次元容量一致条件でTM RMSE \(0.001342\)、Fourier RMSE \(0.002614\)が得られている。ただしこれは再現すべき正解値ではなく、データ生成、seed、実装を監査するためのベンチマークである。結果を近づけるためにtestを見て調整してはならない。
\end{reportbox}

\section{E1B-BOOLE-TWO-SOURCE-IDENTIFIABILITY}

\subsection{目的}

高性能な特徴量でも、観測が異なる源を同一のデータへ写す場合は復元不能である。この境界を、二源混合のfull-rank正対照とrank-one負対照で確認する。

独立な二源を
\begin{equation}
  s_{j,t+1}=F_{\alpha_j}(s_{j,t}),
  \qquad j\in\{1,2\}
  \label{eq:two_source_dynamics}
\end{equation}
とし、観測を
\begin{equation}
  \vect y_t=\mat A\vect s_t,
  \qquad
  \vect s_t=(s_{1,t},s_{2,t})^\top
  \label{eq:linear_mixing}
\end{equation}
とする。ターゲットは順序付き\((\alpha_1,\alpha_2)\)である。

\subsection{二つの観測条件}

\subsubsection*{full-rank正対照}
\(\det\mat A\neq0\)の既知行列を用いる。oracle条件では\(\mat A^{-1}\vect y_t\)から各源を復元して特徴を作り、direct条件では混合チャネルから直接特徴を作る。

\subsubsection*{rank-one負対照}
例えば
\begin{equation}
  y_t=s_{1,t}+s_{2,t}
  \label{eq:rank_one_observation}
\end{equation}
だけを観測する。源を入れ替えた\((\alpha_1,\alpha_2)\)と\((\alpha_2,\alpha_1)\)は同じ観測分布を持ち、順序付きターゲットを区別できない。

\subsection{衝突から導く誤差下限}

同じ観測\(y\)に二つのターゲット
\begin{equation}
  \vect a=(a,b),\qquad \vect b=(b,a)
\end{equation}
が対応し、同確率で現れるとする。二乗誤差を最小にする予測は条件付き平均
\begin{equation}
  \widehat{\vect t}(y)
  =\frac{\vect a+\vect b}{2}
  =\left(\frac{a+b}{2},\frac{a+b}{2}\right)
  \label{eq:collision_optimal_predictor}
\end{equation}
である。各座標の二乗誤差は\((a-b)^2/4\)なので、座標平均RMSEの理論下限は
\begin{equation}
  \mathrm{RMSE}_{\min}
  =\frac{|a-b|}{2}.
  \label{eq:collision_rmse_lower_bound}
\end{equation}
複数ペアでは\(\sqrt{\E[(a-b)^2]/4}\)が対応する。

\begin{theorybox}
\cref{eq:collision_rmse_lower_bound}は、TM、Fourier、ニューラルネットなど手法に依存しない。観測写像が順序を捨てたために生じる情報論的下限である。特徴量を増やしても越えられない。
\end{theorybox}

\subsection{実験設計}

\begin{enumerate}
  \item \(\alpha_1\neq\alpha_2\)のペアでは、同じ二本の源軌道を順序だけ入れ替えたswap標本を作り、元標本とswap標本を必ず同じsplitへ入れる。これによりrank-one観測は標本ごとに厳密に一致する。
  \item full-rank oracle、full-rank direct、rank-one directを比較する。
  \item TMとFourierの特徴容量を一致させる。
  \item rank-oneで、swapペアの特徴距離を直接測る。
  \item 予測RMSEと\cref{eq:collision_rmse_lower_bound}を比較する。
\end{enumerate}

\begin{gatebox}
\begin{enumerate}
  \item full-rank oracleが低誤差であること。これがpipelineの正対照となる。
  \item rank-oneのswap特徴距離が数値誤差内で0に近いこと。
  \item rank-one誤差が理論下限を不自然に下回らないこと。下回れば漏洩を疑う。
  \item rank-oneで低誤差が出た場合、target順序やsample IDが特徴へ漏れていないか確認する。
\end{enumerate}
\end{gatebox}

\begin{reportbox}
既存runでは、full-rank oracle TM RMSE \(0.002083\)、rank-one理論下限 \(0.073030\)、rank-one TM RMSE \(0.073710\)が報告され、swap特徴衝突が確認された。結論は「TMが失われた順序を回復した」ではなく、\textbf{回復不能性を理論下限まで再現した}ことである。
\end{reportbox}

\section{E1の考察課題}

\begin{taskbox}
\textbf{課題 4.1}\quad TMの\(k\)、lag、Fourier帯域数を変え、特徴次元と性能の関係を描け。最良値だけでなく、探索自由度を報告せよ。

\medskip
\textbf{課題 4.2}\quad rank-one観測のターゲットを順序なし集合\(\{\alpha_1,\alpha_2\}\)へ変えると、識別可能性はどう変わるか。適切な損失関数を定式化せよ。

\medskip
\textbf{課題 4.3}\quad full-rankだが条件数の大きい\(\mat A\)を用い、可逆性と数値安定性の違いを調べよ。
\end{taskbox}

% ============================================================
\chapter{E2: 観測雑音・欠測・有限長への頑健性}
\label{chap:e2}
% ============================================================

\begin{purposebox}
E1で得た識別性が、現実的な観測劣化の下でも保たれるかを調べる。clean条件で学習した前処理・特徴・decoderを固定し、test時だけ観測条件を変える。
\end{purposebox}

\section{E2A-BOOLE-OBSERVATION-ROBUSTNESS}

\subsection{研究質問}

\begin{quote}
二源Boole軌道のパラメータ読み出しは、Gaussian雑音、欠測、短い観測窓に対してどのように劣化するか。TMとFourierの差は、基底の性質によるのか、それとも欠測処理の違いによるのか。
\end{quote}

E2AはE1Bのfull-rank観測を引き継ぐ。独立源\(\vect s_t\)を軌道内でrobust標準化し、既知の可逆行列\(\mat A\)で
\begin{equation}
  \vect y_t=\mat A\vect s_t
  \label{eq:e2_clean_observation}
\end{equation}
と観測する。

\subsection{clean-fit契約}

次をclean train/validationだけでfitし、その後は固定する。

\begin{itemize}
  \item median、half-IQR等の前処理規則。
  \item TM次数・lag、Fourier帯域分割。
  \item oracle/directの読み出し経路。
  \item ridge正則化係数。
  \item 欠測補完規則。
\end{itemize}

testで雑音強度や欠測率を見て再fitしてはならない。

\section{劣化モデル}

\subsection{加法Gaussian雑音}

clean観測に対して
\begin{equation}
  \widetilde{\vect y}_t
  =\vect y_t+\sigma\vect\epsilon_t,
  \qquad
  \vect\epsilon_t\sim\mathcal N(\vect0,\mat I)
  \label{eq:e2_gaussian_noise}
\end{equation}
を加える。\(\sigma\in\{0,0.05,0.1,0.2,0.3\}\)を推奨する。標準化前後のどちらで\(\sigma\)を定義したかを明記する。

\subsection{ランダム欠測}

各チャネル・時刻を独立に確率\(r\)で欠測させる。
\begin{equation}
  m_{c,t}\sim\mathrm{Bernoulli}(1-r),
  \qquad r\in\{0,0.1,0.2,0.4,0.6\}.
  \label{eq:e2_missing_mask}
\end{equation}

TM相関では、\(t\)と\(t+\ell\)の両方が観測されたpairだけを平均する。
\begin{equation}
  \widehat C_k^{\mathrm{obs}}(\ell)
  =\frac{\sum_t m_t m_{t+\ell}\overline{M_k(t)}M_k(t+\ell)}
  {\sum_t m_tm_{t+\ell}}.
  \label{eq:pairwise_complete_tm}
\end{equation}
分母が最小pair数を下回る場合は欠損特徴とし、indicatorを追加するか、その条件を失敗として記録する。

Fourierは等間隔系列を要求するため、線形補間、zero-fill、Lomb--Scargle等の選択が必要である。E2Aの基本条件では線形補間を使ってよいが、比較には補間仮定が含まれる。

\subsection{有限観測長}

元の長さ\(T_{\max}\)から先頭またはランダム区間を切り出し、
\begin{equation}
  T\in\{64,128,256,512,1024\}
  \label{eq:e2_observation_lengths}
\end{equation}
を比較する。同じ元軌道の異なる長さをpairedに評価すると分散を減らせる。

\subsection{確認拡張}

E2Bとして次を追加する。

\begin{itemize}
  \item 外れ値を含むimpulsive noise。
  \item Cauchy観測雑音。
  \item 連続block欠測。
  \item 混合行列\(\mat A\)未知のblind条件。
  \item 時刻揺らぎ、不等間隔サンプリング。
\end{itemize}

\section{比較モデル}

\begin{table}[htbp]
\centering
\caption{E2Aの四つの主要経路。各チャネル16次元とし、合計32次元へそろえる。}
\label{tab:e2_models}
\begin{tabularx}{\textwidth}{p{36mm}p{31mm}X}
\toprule
モデル & 観測経路 & 特徴 \\
\midrule
oracle-TM & \(\mat A^{-1}\widetilde{\vect y}\) & \(k=1,\dots,4\)、lag 1,2の複素TM相関 \\
oracle-Fourier & \(\mat A^{-1}\widetilde{\vect y}\) & 16帯域log-energy \\
direct-TM & 混合チャネルを直接使用 & 同じTM相関 \\
direct-Fourier & 混合チャネルを直接使用 & 同じFourier特徴 \\
\bottomrule
\end{tabularx}
\end{table}

oracleは源分離が既知の上限対照、directは観測空間での実用条件である。oracleが常に強いとは限らない。\(\mat A^{-1}\)が雑音を増幅する場合、directが頑健になることもある。

\section{評価と統計}

各劣化レベル、モデル、seedについてRMSEを保存する。cleanからの相対劣化を
\begin{equation}
  D(\eta)
  =\frac{\mathrm{RMSE}(\eta)-\mathrm{RMSE}(0)}
  {\max(\mathrm{RMSE}(0),\epsilon)}
  \label{eq:relative_degradation}
\end{equation}
で定義する。ここで\(\eta\)は\(\sigma\)、\(r\)、または\(T^{-1}\)を表す。

TMとFourierの差は同じsample・seedでpaired比較する。平均だけでなく、中央値、95\%bootstrap区間、最悪seedを報告する。

\begin{gatebox}
\begin{enumerate}
  \item clean性能を再現し、劣化を加える前からpipelineが壊れていない。
  \item 中程度の条件\((\sigma=0.1,r=0.2,T=256)\)で、性能がchanceまたは定数予測まで崩壊しない。
  \item TM優位を主張する場合、欠測処理を入れ替えた感度分析でも結論が保たれる。
  \item oracleの悪化が大きい場合、\(\kappa(\mat A)\)と雑音増幅を報告する。
  \item すべての条件で有効pair数と非有限特徴数を保存する。
\end{enumerate}
\end{gatebox}

\begin{reportbox}
既存runでは、欠測時にTMがFourierを大きく上回る条件があった一方、Gaussian雑音ではFourierが優位な条件もあった。これは「TMが普遍的に頑健」という結論ではない。TMはpairwise-complete平均、Fourierは線形補間という異なる欠測仮定を用いたため、基底と前処理の効果を分離する追加実験が必要である。
\end{reportbox}

\section{可視化}

必須図は次の三つである。

\begin{enumerate}
  \item 横軸を劣化レベル、縦軸をtest RMSEとした曲線。seed区間を付ける。
  \item cleanに対する相対劣化\(D(\eta)\)の曲線。
  \item sampleごとのTM誤差対Fourier誤差の散布図。対角線を入れる。
\end{enumerate}

欠測では、有効pair数分布も図または表で示す。

\begin{taskbox}
\textbf{課題 5.1}\quad 同じ欠測maskに対し、TMにも線形補間、Fourierにもpairwise covarianceを使うなど、前処理を交換した2\(\times\)2実験を設計せよ。

\medskip
\textbf{課題 5.2}\quad \(\mat A\)の条件数を操作し、oracle逆混合が雑音を増幅する量を特異値分解から予測せよ。

\medskip
\textbf{課題 5.3}\quad iid欠測とblock欠測で、同じ欠測率でも時間相関推定の分散が異なる理由を説明せよ。
\end{taskbox}

% ============================================================
\chapter{E3: 同期転移・大域basin・入力保持}
\label{chap:e3}
% ============================================================

\begin{purposebox}
同期を一つの数値だけで判定せず、\textbf{同期多様体の不変性、局所横安定性、有限差初期値からの大域到達、入力情報の保持}を分離して調べる。本章は研究の最小中核である。
\end{purposebox}

\section{同期を四つに分解する}

ノード状態を\(\vect x_t\in\R^N\)、平均状態を
\begin{equation}
  \bar x_t=\frac1N\sum_{i=1}^{N}x_{i,t}
\end{equation}
とする。完全同期多様体は
\begin{equation}
  \mathcal S=\{\vect x:s\ones,\ s\in\R\}
  \label{eq:synchronization_manifold}
\end{equation}
である。

\subsubsection*{不変性}
\(\vect x_t\in\mathcal S\)なら\(\vect x_{t+1}\in\mathcal S\)となる性質である。不変でない多様体に対して横Lyapunov指数を計算しても、完全同期の安定性を表さない。

\subsubsection*{局所横安定性}
\(\mathcal S\)の近傍の微小摂動が平均的に縮むこと、すなわち\(\lambda_\perp<0\)である。

\subsubsection*{大域basin}
有限差を持つ初期値が、極や別アトラクタを越えて実際に同期へ到達する割合である。\(\lambda_\perp<0\)だけでは大域到達を保証しない。

\subsubsection*{情報保持}
同期後または過渡軌道から入力\(u\)を識別・復元できることである。\(E_{\mathrm{sync}}\to0\)でも入力表現がcollapseすれば失敗である。

\section{同期指標}

通常のRMS同期誤差を
\begin{equation}
  E_{\mathrm{sync}}(t)
  =\sqrt{\frac1N\sum_{i=1}^{N}(x_{i,t}-\bar x_t)^2}
  \label{eq:sync_rms}
\end{equation}
とする。Cauchy裾の影響を見るため、robust指標
\begin{equation}
  E_{\mathrm{med}}(t)
  =\median_i |x_{i,t}-\median_jx_{j,t}|
  \label{eq:sync_robust}
\end{equation}
も併記する。

同期成功は一時刻の閾値通過ではなく、終端窓\(\mathcal W\)で
\begin{equation}
  \max_{t\in\mathcal W}E_{\mathrm{sync}}(t)<\varepsilon_{\mathrm{sync}}
  \label{eq:sustained_sync}
\end{equation}
と定義する。推奨は\(\varepsilon_{\mathrm{sync}}=10^{-8}\)だが、float32では別基準を事前に定める。

\section{E3A: 無限次元ランダム結合原稿の監査}

\subsection{原稿のモデルと主張}

添付原稿は、一般化Boole写像を持つ加法的ランダム結合系
\begin{equation}
  x_i(t+1)
  =F_\alpha(x_i(t))
   +\frac{K}{N-1}\sum_{j\ne i}\varepsilon_{ij}x_j(t)
  \label{eq:source_random_additive_gcm}
\end{equation}
を考え、Cauchy尺度の平均場方程式
\begin{equation}
  \gamma_{t+1}
  =\alpha\left(\gamma_t+\frac1{\gamma_t}\right)
   +Km\gamma_t,
  \qquad m=\E|\varepsilon_{ij}|
  \label{eq:source_scale_meanfield}
\end{equation}
を主張する\cite{infinite_sync}。固定点は
\begin{equation}
  \gamma_*^2
  =\frac{\alpha}{1-\alpha-Km}
  \label{eq:source_random_gamma}
\end{equation}
であり、実数正の尺度が存在するためには\(1-\alpha-Km>0\)が必要である。

さらに条件付きLyapunov指数と臨界結合を
\begin{align}
  \lambda_c^{\mathrm{src}}
  &=2\log\left[
    \sqrt\alpha+\sqrt{1-\alpha-Km}
  \right],
  \label{eq:source_conditional_lyap}\\
  K_c^{\mathrm{src}}
  &=\frac{2\sqrt\alpha(1-\sqrt\alpha)}{m}
  \label{eq:source_kc}
\end{align}
とする。

\subsection{最初に行う不変性監査}

\SourceClaim \cref{eq:source_random_additive_gcm}で全ノードを\(x_i=s\)と置くと
\begin{equation}
  x_i(t+1)
  =F_\alpha(s)
   +\frac{Ks}{N-1}\sum_{j\ne i}\varepsilon_{ij}.
  \label{eq:row_sum_sync_update}
\end{equation}
したがって、完全同期多様体が厳密に不変であるには、行和
\begin{equation}
  r_i=\frac1{N-1}\sum_{j\ne i}\varepsilon_{ij}
  \label{eq:row_mean}
\end{equation}
が全\(i\)で等しい必要がある。一般のi.i.d.有限\(N\)行列では、行和揺らぎが残る。

\begin{auditbox}
E3Aの最初の課題は、臨界値を走査することではなく、\cref{eq:row_sum_sync_update}をコードと数式で確認することである。無限次元で行和差が消えるという議論と、有限\(N\)で厳密な同期多様体が存在することは同じではない。
\end{auditbox}

\subsection{三つの実験プロトコル}

\begin{table}[htbp]
\centering
\caption{E3Aで混同してはならない三つのモデル。}
\label{tab:e3a_protocols}
\begin{tabularx}{\textwidth}{p{24mm}p{45mm}X}
\toprule
Protocol & モデル & 目的 \\
\midrule
A & 原稿どおりの未正規化対称ランダム結合 & source claimを直接再現し、行和揺らぎと擬似同期を測る。\\
B & 行和をそろえた対称結合またはregular weighted graph & 同期多様体を保った上で、結合分布の効果を比較する。\\
C & 次節のrow-stochastic output-mixing & 同期理論と入力保持を明確に分離する主モデル。\\
\bottomrule
\end{tabularx}
\end{table}

Protocol AとCの臨界式は異なる。特に\(\alpha=1/2,m=1\)では\cref{eq:source_kc}は約\(0.4142\)であるが、後述するoutput-mixing主モデルは\(K_c=0.5\)となる。同じ記号\(K\)でもモデルが違えば比較できない。

\subsection{同じ\(m=\E|\varepsilon|\)を持つ分布比較}

Protocol AまたはBでは、最低限、次を比較する。

\begin{enumerate}
  \item 一様正結合: \(\varepsilon=1\)。
  \item 正の二点分布: \(\varepsilon\in\{0,2\}\)を同確率。
  \item biased sign: \(P(\varepsilon=1)=0.875\)、\(P(\varepsilon=-1)=0.125\)。
  \item Rademacher: \(P(\varepsilon=\pm1)=1/2\)。
  \item 必要に応じ、\(\E|\varepsilon|=1\)へ正規化したGaussian。
\end{enumerate}

各行列について、\(m\)、\(\mu=\E\varepsilon\)、行和標準偏差
\begin{equation}
  \sigma_r(N)
  =\operatorname{std}_i r_i
  \label{eq:row_sum_std}
\end{equation}
を保存する。経験的閾値が\(m\)と\(\mu\)のどちらに整合するか、また\(\sigma_r(N)\)と同期誤差がどう関係するかを見る。

\begin{gatebox}
\begin{enumerate}
  \item 同期多様体の不変性残差を、各protocolで明示する。
  \item source式との一致を主張する場合、使用したモデルが\cref{eq:source_random_additive_gcm}と一致する。
  \item \(N\)とseedを増やし、経験的閾値の区間を報告する。
  \item \(m\)だけの普遍性を主張する前に、\(\mu\)、行和揺らぎ、符号構造を対照する。
  \item 不一致を数値clipや行列再正規化で隠さない。再正規化は別protocolとする。
\end{enumerate}
\end{gatebox}

\section{E3B-TANGENT-TWO-NODE-BASIN}

\subsection{局所写像}

生のtangent写像を
\begin{equation}
  T_\beta(x)=\tan(\beta x),
  \qquad \beta>1
  \label{eq:raw_tangent_map}
\end{equation}
とする。極は
\begin{equation}
  x=\frac{\pi/2+k\pi}{\beta},
  \qquad k\in\mathbb Z
  \label{eq:tangent_poles}
\end{equation}
にある。

\Derived Cauchy型のポール力学を仮定すると、尺度は
\begin{equation}
  \gamma_{t+1}=\tanh(\beta\gamma_t)
  \label{eq:tangent_scale_map}
\end{equation}
で更新され、\(\beta>1\)では正固定点
\begin{equation}
  \gamma_*=\tanh(\beta\gamma_*)
  \label{eq:tangent_gamma_fixed}
\end{equation}
を持つ。Lyapunov指数の候補閉形式は
\begin{equation}
  \lambda_T(\beta)
  =\log\beta+2\log(1+\gamma_*)
  \label{eq:tangent_lyap}
\end{equation}
である。これはE3Bの前に単一軌道で数値照合する。

\subsection{出力交差混合型}

\begin{align}
  x_{t+1}&=(1-\eta)T_\beta(x_t)+\eta T_\beta(y_t),\\
  y_{t+1}&=\eta T_\beta(x_t)+(1-\eta)T_\beta(y_t).
  \label{eq:output_cross_tangent}
\end{align}
同期軌道\(x_t=y_t=s_t\)では\(s_{t+1}=T_\beta(s_t)\)である。微小差\(\delta_t=x_t-y_t\)は
\begin{equation}
  \delta_{t+1}
  \approx(1-2\eta)T'_\beta(s_t)\delta_t
  \label{eq:output_cross_variational}
\end{equation}
に従うため、横Lyapunov指数は
\begin{equation}
  \lambda_\perp^{\mathrm{out}}
  =\lambda_T(\beta)+\log|1-2\eta|.
  \label{eq:output_cross_transverse}
\end{equation}
局所同期条件は
\begin{equation}
  |1-2\eta|<\e^{-\lambda_T(\beta)}.
  \label{eq:output_cross_interval}
\end{equation}
特に\(\eta=1/2\)では一回の更新で\(x_{t+1}=y_{t+1}\)となる。これは同期判定コードの正対照である。

\subsection{状態拡散型}

\begin{align}
  x_{t+1}&=T_\beta(x_t)+\kappa(y_t-x_t),\\
  y_{t+1}&=T_\beta(y_t)+\kappa(x_t-y_t).
  \label{eq:state_diffusive_tangent}
\end{align}
同期軌道は同じく\(s_{t+1}=T_\beta(s_t)\)であり、横方向は
\begin{equation}
  \delta_{t+1}
  \approx\left[T'_\beta(s_t)-2\kappa\right]\delta_t
  \label{eq:state_diffusive_variational}
\end{equation}
となる。したがって
\begin{equation}
  \lambda_\perp^{\mathrm{diff}}
  =\lim_{T\to\infty}\frac1T
    \sum_{t=0}^{T-1}
    \log\left|\beta\sec^2(\beta s_t)-2\kappa\right|
  \label{eq:state_diffusive_transverse}
\end{equation}
を同期軌道から推定する。

\subsection{局所と大域の二段評価}

\begin{enumerate}
  \item \textbf{局所}: 同期軌道\(s_t\)を4万点以上生成し、\cref{eq:output_cross_transverse,eq:state_diffusive_transverse}を計算する。
  \item \textbf{大域}: 独立なCauchy初期対\((x_0,y_0)\)を最低256組生成し、500ステップ反復する。
  \item 終端50ステップすべてで\(|x_t-y_t|<10^{-8}\)なら成功とする。
  \item 発散率、同期到達時間、一時最大誤差、同期後再離脱回数を保存する。
\end{enumerate}

推奨走査は
\begin{equation}
  \beta\in\{1.05,1.10,1.20,1.30\}
  \label{eq:e3b_beta_grid}
\end{equation}
であり、\(\eta\)と\(\kappa\)は粗走査後、\(\lambda_\perp=0\)近傍を細かくする。

\begin{cautionbox}
GPUの高速\code{tan}、CPUのlibm、float32、float64では極近傍の引数還元が異なり得る。暗黙のmoduloやclipを入れず、計算精度とバックエンドごとにrunを分ける。発散した軌道を成功率の分母から除外してはならない。
\end{cautionbox}

\begin{gatebox}
\begin{enumerate}
  \item 出力混合の理論同期区間\cref{eq:output_cross_interval}を局所指数で再現する。
  \item \(\eta=1/2\)の一ステップ同期が全seedで通る。
  \item 状態拡散型で\(\lambda_\perp<0\)の領域と大域成功率を別図にする。
  \item 局所安定だが大域失敗する領域を、発散率と極近傍通過で説明する。
  \item 精度を変えても主要相図が定性的に保たれる。
\end{enumerate}
\end{gatebox}

\begin{reportbox}
既存runでは、局所的に\(\lambda_\perp<0\)であったgrid点のうち、大域成功率が0.9未満の割合が約0.55であった。これは「局所横安定性だけから任意の有限差初期値の同期を結論できない」という境界を示す。\(\eta=1/2\)は成功率1の正対照だが、情報圧縮モデルとしては強制的すぎる。
\end{reportbox}

\section{E3C-BOOLE-N8-INPUT-RETENTION}

\subsection{同期多様体を保証する主モデル}

以降の情報保持実験では、局所写像の\textbf{出力}をrow-stochastic行列で混合する。主実験ではgraph modeを直交基底として扱えるよう、\(\mat W\)を対称row-stochastic行列に限定する。非対称・非正規行列は固有ベクトルの非直交性を別に扱う拡張実験とする。
\begin{equation}
  \vect x_{t+1}
  =\mat B_K F_\alpha(\vect x_t)
   +g u_t\ones,
  \qquad
  \mat B_K=(1-K)\mat I+K\mat W,
  \label{eq:output_mixing_network}
\end{equation}
ここで\(F_\alpha\)は成分ごとに作用し、\(\mat W\ones=\ones\)とする。全ノードが\(s_t\)なら
\begin{equation}
  \vect x_{t+1}
  =\left(F_\alpha(s_t)+gu_t\right)\ones
\end{equation}
となるため、同期多様体は厳密に不変である。

\subsection{graph modeごとの横指数}

\(\mat W\)の固有値を\(\nu_q\)、固有ベクトルを\(\vect v_q\)とする。同期方向は\(\nu_0=1\)、\(\vect v_0\propto\ones\)である。同期軌道周りのgraph mode\(q\)の摂動\(a_{q,t}\vect v_q\)は
\begin{equation}
  a_{q,t+1}
  \approx F'_\alpha(s_t)
  \left(1-K+K\nu_q\right)a_{q,t}
  \label{eq:graph_mode_variational}
\end{equation}
に従う。よって
\begin{equation}
  \lambda_q
  =\lambda_\parallel
   +\log|1-K+K\nu_q|
  \label{eq:graph_mode_lyapunov}
\end{equation}
である。共通入力\(gu_t\ones\)は横方向差で相殺されるが、同期軌道\(s_t\)を変えるため、\(\lambda_\parallel\)は実際の駆動軌道から測る。

\subsubsection*{全平均混合}
\(\mat W=\ones\ones^\top/N\)なら、すべての横モードで\(\nu_q=0\)である。自律系では
\begin{equation}
  \lambda_\perp
  =\lambda_F+\log|1-K|.
  \label{eq:meanfield_output_transverse}
\end{equation}
\(0\le K\le1\)での下側臨界は
\begin{equation}
  K_c^{\mathrm{out}}
  =1-\e^{-\lambda_F}.
  \label{eq:output_mixing_kc}
\end{equation}
特に\(\alpha=1/2\)では\(\lambda_F=\log2\)なので
\begin{equation}
  K_c^{\mathrm{out}}=0.5.
  \label{eq:output_mixing_kc_half}
\end{equation}

\begin{cautionbox}
\cref{eq:output_mixing_kc}は\cref{eq:source_kc}とは別の式である。前者はrow-stochasticな出力混合、後者は添付原稿の加法的ランダム結合に対するsource claimである。モデル名を省略して「理論臨界値」とだけ書かない。
\end{cautionbox}

\subsection{入力の埋込み}

E3Cでは\(N=8\)、\(\alpha=1/2\)とし、二値入力\(u\in\{-1,+1\}\)を初期横モードへ埋め込む。第一群と第二群に反対符号を持つ単位ベクトル
\begin{equation}
  v_i=
  \begin{cases}
    +1/\sqrt N,&i\le N/2,\\
    -1/\sqrt N,&i>N/2
  \end{cases},
  \qquad \ones^\top\vect v=0
  \label{eq:contrast_graph_mode}
\end{equation}
を作り、
\begin{equation}
  \vect x_0
  =c\ones+a u\vect v+\vect\xi
  \label{eq:e3c_initial_encoding}
\end{equation}
とする。\(c\)は共通Cauchy成分、\(a\)は小さい入力振幅、\(\vect\xi\)はzero-sum nuisanceである。E3Cでは入力を初期条件だけへ入れ、\(t\ge0\)の共通外場は\(g u_t=0\)とする。同じ\((c,\vect\xi)\)に対する\(u=-1,+1\)をpaired sampleとし、pair単位でsplitする。

\subsection{TM--graph readout}

終端\(L\)ステップの各ノードにTMモード
\begin{equation}
  M_{k,i}(t)
  =\left(\frac{x_{i,t}-\ii\gamma_*}{x_{i,t}+\ii\gamma_*}\right)^k,
  \qquad k=1,\dots,4
  \label{eq:e3c_node_tm}
\end{equation}
を適用する。uniform modeとcontrast modeへ
\begin{align}
  z_{k,0}(t)&=\frac1{\sqrt N}\sum_iM_{k,i}(t),\\
  z_{k,1}(t)&=\sum_i v_iM_{k,i}(t)
  \label{eq:e3c_graph_projection}
\end{align}
と射影し、それぞれの時間平均とlag-1複素自己相関の実部・虚部を連結する。\(k=1,\dots,4\)、graph mode 2種、統計4実数で合計32次元である。

比較は
\begin{itemize}
  \item TM--graph 32次元。
  \item raw graph時間統計 32次元。
  \item terminal state 8次元。
\end{itemize}
とする。decoderは同じlogistic regressionを用い、balanced accuracyとlog lossを報告する。

\subsection{同期と情報のPareto評価}

\(K\)ごとに
\begin{equation}
  \left(
    \median E_{\mathrm{sync}},
    \mathrm{BA}_{\mathrm{input}}
  \right)
  \label{eq:sync_info_pair}
\end{equation}
を描く。左上、すなわち低同期誤差・高識別性能が望ましい。

\begin{figure}[htbp]
\centering
\begin{tikzpicture}[x=1cm,y=1cm]
\draw[-{Latex[length=2.5mm]},line width=0.8pt] (0,0)--(10.2,0) node[right]{\small 同期誤差};
\draw[-{Latex[length=2.5mm]},line width=0.8pt] (0,0)--(0,6.2) node[above]{\small 入力識別性能};
\fill[LightTeal] (0,4.6) rectangle (3.0,6.0);
\node[Forest,align=center] at (1.5,5.3) {有用な\\同期縮約候補};
\fill[LightRed] (0,0) rectangle (3.0,1.6);
\node[Brick,align=center] at (1.5,0.8) {collapse\\低情報};
\fill[LightOrange] (7.0,4.6) rectangle (10.0,6.0);
\node[Orange,align=center] at (8.5,5.3) {高情報だが\\非同期reservoir};
\draw[dashed,MidGray] (3,0)--(3,6);
\draw[dashed,MidGray] (0,4.6)--(10,4.6);
\fill[DeepNavy] (4.8,3.9) circle (2.4pt) node[above right]{\small 臨界近傍候補};
\fill[DeepNavy] (1.0,1.1) circle (2.4pt) node[above right]{\small 深い同期};
\fill[DeepNavy] (8.1,5.1) circle (2.4pt) node[below left]{\small 弱結合};
\end{tikzpicture}
\caption{同期誤差と入力情報を同じスカラーへ混ぜず、Pareto平面として評価する。}
\label{fig:sync_information_pareto}
\end{figure}

\subsection{探索条件}

\begin{equation}
  K\in\{0,0.30,0.45,0.50,0.52,0.60,0.80,0.95\}
  \label{eq:e3c_k_grid}
\end{equation}
を基準とし、\(K=0.5\)近傍を細かくする。終端窓は24ステップを基準とする。入力振幅\(a\)とnuisance尺度はvalidationで少数候補から選び、testを使わない。

\begin{gatebox}
\begin{enumerate}
  \item \(K>K_c^{\mathrm{out}}\)で同期誤差が大きく低下する。
  \item 入力BAの95\%下限が0.5を上回る同期側の点が存在する。
  \item 強い主張にはBA\(\ge0.65\)を要求する。
  \item 深い同期でBAがchanceへ戻るなら、情報collapseとして報告する。
  \item effective rankと同期誤差を併記し、「低rankだから良い」とは結論しない。
\end{enumerate}
\end{gatebox}

\subsection{effective rank}

潜在特徴行列\(\mat Z\)の特異値を\(\sigma_1,\dots,\sigma_r\)とすると、participation ratio型のeffective rankを
\begin{equation}
  \effrank{}(\mat Z)
  =\frac{\left(\sum_j\sigma_j^2\right)^2}
  {\sum_j\sigma_j^4}
  \label{eq:effective_rank}
\end{equation}
とする。値が小さければ分散が少数方向へ集中しているが、その方向が入力を表すとは限らない。

\begin{reportbox}
既存pilotでは、同期側の最良TM条件が\(K=0.52\)、balanced accuracy \(0.59174\)、median sync RMS \(0.00621\)であった。同期誤差低下は支持されたが、事前の強いBA\(\ge0.65\)は未支持であった。また深い同期では情報低下が見られた。この結果は「臨界近傍が常に最適」という証明ではなく、より良い入力注入とreadoutを必要とする反例・候補点である。
\end{reportbox}

\section{E3の考察課題}

\begin{taskbox}
\textbf{課題 6.1}\quad \(\mat W=\ones\ones^\top/N\)と、自己結合を除く\(\mat W=(\ones\ones^\top-\mat I)/(N-1)\)について、横固有値と臨界\(K\)を導け。

\medskip
\textbf{課題 6.2}\quad E3Cの入力を、(i)初期横モード、(ii)共通外場、(iii)結合係数変調の三方式へ入れ、同一budgetで比較するE3Dを設計せよ。

\medskip
\textbf{課題 6.3}\quad 完全同期ではなく2-cluster同期を目標にする結合行列を設計し、どのgraph固有モードを中立または弱収縮に残すべきか説明せよ。

\medskip
\textbf{課題 6.4}\quad E3AのProtocol Aで、行和揺らぎ\(\sigma_r(N)\)が\(N^{-1/2}\)程度に減るかを調べ、同期誤差との関係を回帰せよ。
\end{taskbox}

% ============================================================
\chapter{E4: 合成時系列の圧縮・復元}
\label{chap:e4}
% ============================================================

\begin{purposebox}
同期系から得た低次元表現が、単なる入力分類ではなく波形そのものを復元できるかを調べる。潜在次元をrate、復元誤差をdistortionとして、TM、raw state、terminal state、入力PCAを公平に比較する。
\end{purposebox}

\section{E4A: 合成信号の復元}

\noindent\textbf{正規実験ID:} \EID{E4A-BOOLE-SYNTHETIC-SIGNAL-RECONSTRUCTION}

\subsection{研究質問}

\begin{quote}
共通入力で駆動された同期Booleネットワークの有限時間軌道は、元信号を少数次元から復元できるか。TM時間表現は、同じ潜在次元・同じridge decoderで、raw time featureや終端状態を上回るか。
\end{quote}

\section{合成信号}

長さを\(L=64\)、正規化時刻を\(\tau_n=n/(L-1)\)とする。少なくとも次の三族を生成する。

\subsection{正弦波}
\begin{equation}
  u_n=A\sin(2\pi f\tau_n+\phi).
  \label{eq:e4_sine}
\end{equation}
推奨範囲は\(A\in[0.5,1.5]\)、\(f\in[1,8]\) cycles/window、\(\phi\in[0,2\pi)\)である。

\subsection{chirp}
\begin{equation}
  u_n=A\sin\left[2\pi\left(f_0\tau_n+\frac12r\tau_n^2\right)+\phi\right].
  \label{eq:e4_chirp}
\end{equation}
\(f_0\)と\(r\)を独立に変える。最終瞬時周波数がNyquistを越えない範囲に制約する。

\subsection{二周波混合}
\begin{equation}
  u_n=A_1\sin(2\pi f_1\tau_n+\phi_1)
      +A_2\sin(2\pi f_2\tau_n+\phi_2).
  \label{eq:e4_two_tone}
\end{equation}
\(f_1\)と\(f_2\)の差が小さい条件と大きい条件を含める。

各信号は最大絶対値またはtrain全体の尺度で正規化する。sampleごとの振幅正規化を行うと振幅情報が消えるため、何をターゲットとして残すかに応じて選ぶ。

\section{駆動同期ネットワーク}

E3Cと同じ出力混合系
\begin{equation}
  \vect x_{n+1}
  =\mat B_KF_\alpha(\vect x_n)+g u_n\ones
  \label{eq:e4_driven_network}
\end{equation}
を用いる。基準条件は\(N=8\)、\(\alpha=1/2\)である。

共通入力はノード差から直接消えるが、同期軌道\(s_n\)を変える。したがって、自律系の\(K_c=0.5\)をそのまま保証として使わず、各\((K,g)\)で経験的横指数
\begin{equation}
  \widehat\lambda_\perp(K,g)
  =\frac1L\sum_{n=0}^{L-1}
    \log\left|F'_\alpha(s_n)(1-K)\right|
  \label{eq:e4_driven_transverse}
\end{equation}
と同期誤差を測る。

\begin{cautionbox}
共通入力が横方向で相殺されることと、入力が同期安定性に影響しないことは同じではない。\(F'(s_n)\)が入力で変わるため、平均収縮率も変わり得る。
\end{cautionbox}

\section{データ分割}

train/validation/testは信号族で層化し、元パラメータseed単位で分ける。より厳しいOOD試験として、testだけ周波数帯、chirp rate、振幅範囲をずらす。

力学パラメータ\((K,g)\)は、validationのTM表現・潜在16次元のNMSEだけで選ぶ。testを使わない。画像へ進むE5Aでは、この\((K,g)\)を凍結する。

\section{比較表現}

\begin{table}[htbp]
\centering
\caption{E4Aで比較する表現。PCAとdecoderはtrainだけでfitする。}
\label{tab:e4_representations}
\begin{tabularx}{\textwidth}{p{35mm}X}
\toprule
表現 & 定義と役割 \\
\midrule
TM spectrum & ノード平均Cayley mode \(k=1,\dots,4\)の時間系列を作り、DCTまたは時間相関を連結する提案表現。\\
raw mean DCT & ノード平均の生状態をrobust標準化し、DCT係数を使うTMなし時間baseline。\\
reservoir flat & 全ノード\(\times\)全時刻をflattenし、PCAで圧縮する。軌道に残る情報の上限baseline。\\
terminal state & 最終時刻の\(N\)状態。時間履歴なしbaseline。\\
input oracle & 入力波形自身をPCAで圧縮する。力学を通さない線形上限baseline。\\
\bottomrule
\end{tabularx}
\end{table}

TM spectrumの一例として
\begin{equation}
  \bar M_k(n)=\frac1N\sum_iM_{k,i}(n)
  \label{eq:e4_mean_tm_mode}
\end{equation}
を作り、\(\Re\bar M_k\)、\(\Im\bar M_k\)のDCT係数を連結する。どの係数を保持するかはtrain/validationだけで固定する。

各表現\(\vect h\)をPCAで
\begin{equation}
  d_z\in\{4,8,16,32\}
  \label{eq:e4_latent_dimensions}
\end{equation}
へ圧縮し、同じridge decoder
\begin{equation}
  \widehat{\vect u}
  =\mat W_d\vect z+\vect b_d
  \label{eq:e4_ridge_decoder}
\end{equation}
で64点を復元する。

\section{指標}

平均二乗誤差を入力分散で正規化した
\begin{equation}
  \NMSE{}
  =\frac{\sum_{n}(u_n-\widehat u_n)^2}
  {\sum_n(u_n-\bar u)^2+\epsilon}
  \label{eq:nmse}
\end{equation}
を主要指標とする。併せてRMSE、相関係数、周波数領域誤差を報告する。

潜在次元をrateのproxyとするが、次元はbit数ではない。確認実験では、各潜在座標を\(b\in\{4,8,16\}\) bitへ量子化し、
\begin{equation}
  R=\frac{d_zb}{L}
  \quad\text{[bits/sample]}
  \label{eq:bit_rate_proxy}
\end{equation}
を横軸にしたrate--distortion曲線も作る。

\section{公平比較の条件}

\begin{enumerate}
  \item 全表現で同じtrain/validation/test splitを使う。
  \item PCA次元とridgeの探索候補をそろえる。
  \item reservoir flatだけ元次元が大きいため、計算量と保存量も併記する。
  \item input oracleは提案法ではなく、入力が元々低次元かを示す対照である。
  \item decoderを表現ごとに大幅に変えない。
\end{enumerate}

\begin{gatebox}
\begin{enumerate}
  \item TMがterminal stateを同じ\(d_z\)で上回る。これは時間履歴の必要性を示す。
  \item TMがraw mean DCTを上回る場合にのみ、TM座標固有の利点を主張する。
  \item TMがreservoir flatを少ない元特徴・計算量で近似できるかを調べる。
  \item input oracleの25\%以内のNMSEを強い成功基準とする。
  \item 入力PCAが圧倒的に良い場合、「合成族が元々低次元で力学が不要」という代替説明を採用する。
\end{enumerate}
\end{gatebox}

\begin{reportbox}
既存pilotの潜在16次元では、input oracleのtest NMSEが\(0.00661\)で最良であり、reservoir flat、TM spectrum、raw mean DCTはいずれも約1.0であった。TMはraw DCTとterminalを一部条件で上回ったが、input PCAには遠く及ばなかった。これは、現在の入力注入またはreadoutが波形情報を保持できていない負の結果である。
\end{reportbox}

\section{失敗時の切り分け}

\begin{table}[htbp]
\centering
\caption{E4でNMSEが高いときの診断。}
\label{tab:e4_failure_diagnosis}
\begin{tabularx}{\textwidth}{p{42mm}X}
\toprule
観測結果 & 有力な原因と次の確認 \\
\midrule
reservoir flatも悪い & 力学へ入力情報が入っていない。入力gain、注入位置、初期状態を見直す。\\
reservoir flatは良いがTMが悪い & 軌道には情報が残るが、現在のTM統計が読めていない。高次mode、delay、cross-node特徴を追加する。\\
TMは良いが同期しない & 表現はreservoirとして有用だが、同期圧縮の仮説は未支持。\\
同期するが全表現が悪い & 入力が横・縦両方向で消え、collapseしている可能性。\\
input oracleだけ良い & 入力族が低次元で、非線形力学を通す利点がない。より複雑な信号または別目的を検討する。\\
\bottomrule
\end{tabularx}
\end{table}

\begin{taskbox}
\textbf{課題 7.1}\quad 分枝記号、TM高次mode、delay embeddingを一つずつ追加するE4B ablationを設計せよ。一度に複数変更しないこと。

\medskip
\textbf{課題 7.2}\quad 共通入力、ノード別入力、初期条件入力について、横方向と縦方向のどちらへ情報が入るかを変分方程式で比較せよ。

\medskip
\textbf{課題 7.3}\quad 復元NMSEだけでなく、周波数、振幅、位相の推定誤差へ分解し、どの属性が失われているかを調べよ。
\end{taskbox}

% ============================================================
\chapter{E5: 画像pilotと確認実験}
\label{chap:e5}
% ============================================================

\begin{purposebox}
E4で固定した力学を小画像へ転移し、波形以外の入力でも潜在表現が機能するかを確認する。pilotと本確認を区別し、8\(\times\)8 digitsからMNISTや一般画像への過度な一般化を防ぐ。
\end{purposebox}

\section{E5A-BOOLE-SMALL-IMAGE-DIGITS}

\subsection{データ}

scikit-learnのdigitsデータセットを用いる。1797枚、8\(\times\)8、10クラスである。画素を\([0,1]\)へ正規化し、基本条件ではrow-majorで長さ64の系列
\begin{equation}
  u_0,u_1,\dots,u_{63}
  \label{eq:digits_sequence}
\end{equation}
へ変換する。

train/validation/testはクラス層化し、元画像単位で分ける。複数seedで分割し、同じ画像が別splitへ入らないようにする。

\subsection{転移契約}

E4Aで選んだ\(K\)とinput gain\(g\)を凍結し、画像testに合わせて再調整しない。再学習してよいのは、各表現のtrain-only PCA、ridge decoder、logistic probeである。この契約により、画像を見て力学を選ぶ楽観バイアスを避ける。

\section{二つの評価課題}

\subsubsection*{再構成}
潜在\(\vect z\)から64画素をridgeで復元し、\NMSE{}、\PSNR{}、\SSIM{}を測る。画素最大値を1とすると
\begin{equation}
  \PSNR{}
  =10\log_{10}\frac{1}{\mathrm{MSE}}
  \label{eq:psnr}
\end{equation}
である。

\subsubsection*{クラス保持}
同じ潜在\(\vect z\)からlogistic regressionで数字クラスを予測し、balanced accuracy、macro-F1、log lossを測る。再構成と分類は別decoderで評価する。

\begin{cautionbox}
分類精度が高くても画像が復元できるとは限らず、復元が良くてもクラス境界が線形に読めるとは限らない。二つを別の問いとして報告する。
\end{cautionbox}

\section{比較表現}

E4と同じTM spectrum、raw mean DCT、reservoir flat、terminal state、input oracleを使う。潜在次元も\(d_z\in\{4,8,16,32\}\)へそろえる。画像用に新しい特殊特徴を追加する場合はE5Cとして分離する。

\section{入力走査順の監査}

row-majorは、二次元近傍の一部しか保存しない。次の走査を追加し、力学以前の表現差を調べる。

\begin{enumerate}
  \item row-major。
  \item column-major。
  \item serpentine scan。
  \item Hilbert curveまたはZ-order。
  \item 8系列のrow並列入力。
\end{enumerate}

走査順をvalidationで選ぶ場合、選択自由度として明示する。E5Aの基本再現はrow-majorに固定する。

\begin{gatebox}
\begin{enumerate}
  \item TMがterminal stateを共通\(d_z\)で上回る。
  \item TMがraw DCTを上回る場合にのみTM固有の画像利点を主張する。
  \item class balanced accuracy \(\ge0.70\)をpilotの強い成功基準とする。
  \item input PCAとのSSIM差が0.10以内なら、線形上限への接近とみなす。
  \item どの提案表現もchance分類・NMSE約1なら、画像への転移は未支持とする。
\end{enumerate}
\end{gatebox}

\begin{reportbox}
既存pilotの潜在16次元では、input oracleがNMSE \(0.15326\)、PSNR約\(19.88\)、SSIM \(0.9532\)、balanced accuracy \(0.96124\)で最良であった。TM spectrumはNMSE約\(1.013\)、SSIM約\(0.572\)、balanced accuracy約\(0.102\)で、強いgateは未支持であった。これはE5Aの負の結果であり、MNISTへ進む前に入力注入とreadoutを再設計する根拠となる。
\end{reportbox}

\section{E5B-BOOLE-MNIST-CONFIRMATORY}

E5Bは28\(\times\)28 MNISTを用いる確認実験であり、E5Aとは別EIDにする。E5Aが失敗している状態で、規模だけ増やして実行する優先度は低い。実行する場合は、次を同一budgetで比較する。

\begin{table}[htbp]
\centering
\caption{E5Bで必要なbaseline。}
\label{tab:e5b_baselines}
\begin{tabularx}{\textwidth}{p{36mm}X}
\toprule
モデル & 公平比較条件 \\
\midrule
input PCA & 同じ潜在bit数。線形圧縮上限。\\
linear AE & 同じ潜在次元と線形decoder。PCAとの整合確認。\\
small AE & decoderパラメータ数と学習stepを合わせる。\\
VAE & \(\beta\)、潜在次元、decoder容量、seedを明示。\\
non-chaotic reservoir & ノード数、時間step、readout容量を合わせる。\\
TM-chaos model & 力学パラメータ選択をvalidationだけで行う。\\
\bottomrule
\end{tabularx}
\end{table}

比較では、潜在次元だけでなく、bit数、decoderパラメータ数、学習FLOPs、推論時間、seed数をそろえる。画像圧縮を主張するなら、復元品質に加えて実際のrateを測る。

\section{E5の考察課題}

\begin{taskbox}
\textbf{課題 8.1}\quad row-majorとHilbert走査で、隣接画素が系列上で離れる距離分布を比較せよ。復元差が力学によるものか走査によるものか議論せよ。

\medskip
\textbf{課題 8.2}\quad E5Aの失敗が、(i)入力gain、(ii)共通入力による空間情報消失、(iii)TM統計の過度な時間平均、(iv)線形decoder容量のどれによるか、一要因ずつ切り分ける計画を作れ。

\medskip
\textbf{課題 8.3}\quad 画像を一つの長系列へ潰さず、各rowまたはpatchをノードへ対応させる二次元結合設計を定式化せよ。
\end{taskbox}

% ============================================================
\chapter{先行実装AKOrNとの比較と学習可能な拡張}
\label{chap:akorn}
% ============================================================

\begin{purposebox}
同期をニューラル表現へ組み込んだ先行例AKOrNを、カオス同期エンコーダの設計参考として読む。共通点と相違点を明確にし、単純な模倣ではなく、何を借り、何を独自検証するかを定める。
\end{purposebox}

\section{AKOrNの基本式}

AKOrNは各ニューロンを単位球面上のベクトル\(\vect x_i\in\R^d\)、\(\|\vect x_i\|_2=1\)として表し、一般化Kuramoto型の連続時間更新
\begin{equation}
  \dot{\vect x}_i
  =\mat\Omega_i\vect x_i
   +\Proj_{\vect x_i}\left(
      \vect c_i+\sum_j\mat J_{ij}\vect x_j
    \right)
  \label{eq:akorn_continuous}
\end{equation}
を用いる\cite{akorn}。射影は
\begin{equation}
  \Proj_{\vect x}(\vect y)
  =\vect y-\langle\vect y,\vect x\rangle\vect x
  \label{eq:akorn_projection}
\end{equation}
であり、球面の接空間へ更新方向を移す。\(\mat\Omega_i\)は反対称な自然回転、\(\vect c_i\)は入力依存刺激、\(\mat J_{ij}\)は結合である。

離散時間では
\begin{align}
  \Delta\vect x_i^{(t)}
  &=\mat\Omega_i\vect x_i^{(t)}
   +\Proj_{\vect x_i^{(t)}}\left(
     \vect c_i+\sum_j\mat J_{ij}\vect x_j^{(t)}
   \right),\\
  \vect x_i^{(t+1)}
  &=\Pi\left[\vect x_i^{(t)}+h\Delta\vect x_i^{(t)}\right],
  \qquad \Pi(\vect y)=\frac{\vect y}{\|\vect y\|_2}
  \label{eq:akorn_discrete}
\end{align}
と有限回更新する。readoutでは相対方向を読むため、線形結合のノルム
\begin{equation}
  m_k=\left\|\sum_i\mat U_{ki}\vect x_i^{(T)}\right\|_2
  \label{eq:akorn_readout}
\end{equation}
を用い、global rotationに不変な特徴を作る。

\section{本研究との対応}

\begin{table}[htbp]
\centering
\caption{AKOrNと本研究の対応。}
\label{tab:akorn_comparison}
\begin{tabularx}{\textwidth}{p{31mm}p{53mm}X}
\toprule
観点 & AKOrN & 本研究 \\
\midrule
局所状態 & 球面上の有限次元ベクトル & 実直線上のCauchy型カオス状態 \\
力学 & Kuramoto型回転・同期 & 一般化Booleまたはtangent写像 \\
入力 & symmetry-breaking stimulus \(\vect c_i\) & 初期条件、共通外場、横モード、結合変調 \\
結合 & 学習可能なconv/attentionを含む & まず解析可能な固定結合、後に学習可能化 \\
readout & 相対方向のノルム & Cayley/TM時間相関、graph mode、decoder \\
主評価 & 物体binding、数独、頑健性 & 同期、識別可能性、rate--distortion、復元 \\
理論基盤 & 球面幾何とLyapunov energy & Cauchy安定性、複素極、横Lyapunov \\
\bottomrule
\end{tabularx}
\end{table}

AKOrNは、同期を分散的・連続的クラスタリングとして機械学習へ利用できる先行例である。添付論文では、object discovery、Sudoku、adversarial/common corruptionへの頑健性、calibrationが報告される\cite{akorn}。一方、AKOrNはカオス同期、Cauchy不変測度、TM復号を扱わない。

\section{借りる設計原理}

本研究がAKOrNから借りるべき要素は次である。

\begin{enumerate}
  \item \textbf{有限時間の力学ブロック}: 収束後の一点だけでなく、数stepの過渡を表現として使う。
  \item \textbf{入力刺激と状態の二流}: 入力\(C\)と動的状態\(X\)を分ける。
  \item \textbf{不変量を読むreadout}: global phaseや共通方向に依存しない関係量を読む。
  \item \textbf{ablation}: 射影、正規化、結合、readoutを一つずつ外し、何が性能を生むか調べる。
  \item \textbf{test-time dynamics}: 難しい入力では反復stepを増やしたときの性能変化を見る。
\end{enumerate}

\section{そのまま借りない要素}

\begin{cautionbox}
AKOrNの性能向上を、カオス同期エンコーダの性能証拠として使ってはならない。AKOrNの球面正規化は振幅情報を捨てるため、著者らも事象の存在や長期記憶に弱い可能性を述べている。逆に本研究の実直線カオスは極と発散を持ち、AKOrNの数値安定性をそのまま期待できない。
\end{cautionbox}

\section{学習可能な最小拡張}

固定力学のE0--E5を通過した後、次の形へ拡張する。
\begin{equation}
  \vect x_{t+1}
  =\mat B_{\theta}(u)
    F_{\alpha_\theta(u)}(\vect x_t)
   +\vect c_{\theta}(u),
  \qquad t=0,\dots,T-1
  \label{eq:learnable_chaos_block}
\end{equation}
readoutは
\begin{equation}
  \vect z
  =R_{\phi}\left(
    \{M_k(x_{i,t})\}_{i,k,t}
  \right),
  \qquad
  \widehat{\vect u}=D_\psi(\vect z)
  \label{eq:learnable_tm_encoder}
\end{equation}
とする。

最初から\(\theta,\phi,\psi\)を同時学習すると、どの要素が効いたか分からない。次の順で解凍する。

\begin{enumerate}
  \item 固定力学 + 固定TM + 線形decoder。
  \item 固定力学 + 学習readout + 小decoder。
  \item 学習入力写像 + 固定結合。
  \item 対称性・row-stochastic制約付き学習結合。
  \item 最後に局所パラメータ\(\alpha\)を学習する。
\end{enumerate}

\begin{gatebox}
学習可能化は、固定系の失敗原因が特定されてから行う。学習decoderを大きくしただけで復元できる場合、カオス同期表現の利点とはみなさない。パラメータ数、FLOPs、反復step、メモリ、baselineをそろえる。
\end{gatebox}

\begin{taskbox}
\textbf{課題 9.1}\quad AKOrNの\cref{eq:akorn_readout}と本研究のTM--graph readoutを、「どの群作用に不変か」という観点で比較せよ。

\medskip
\textbf{課題 9.2}\quad \cref{eq:learnable_chaos_block}で同期多様体を保つための\(\mat B_\theta(u)\)の制約を三つ挙げよ。

\medskip
\textbf{課題 9.3}\quad 反復stepをtest時に増やしたとき、復元が改善する場合と情報collapseする場合を予測し、検証曲線を設計せよ。
\end{taskbox}

% ============================================================
\chapter{実行ロードマップと卒論最小スコープ}
\label{chap:roadmap}
% ============================================================

\begin{purposebox}
実装順、時間配分、停止条件を固定する。研究の価値を「すべて成功したこと」ではなく、「どの仮説がどこまで成立するかを切り分けたこと」に置く。
\end{purposebox}

\section{推奨実装順}

\begin{table}[htbp]
\centering
\caption{自分で実装する推奨順。}
\label{tab:implementation_order}
\begin{tabularx}{\textwidth}{p{15mm}p{42mm}X}
\toprule
順 & 実験 & 完了時に言えること \\
\midrule
1 & E0A & Boole写像とCauchy/Lyapunov実装が正しい。\\
2 & E0B & Cayley/TM基底とK-fold正対照が正しい。\\
3 & E1A & 有限軌道の時間順序にパラメータ情報が残るか。\\
4 & E1B & 観測が情報を捨てたときの回復不能境界。\\
5 & E2A & 読み出しの雑音・欠測・有限長頑健性。\\
6 & E3A & 無限次元原稿の有限\(N\)監査と結合分布比較。\\
7 & E3B & 局所横安定性と大域basinの差。\\
8 & E3C & 同期と入力保持のPareto関係。\\
9 & E4A & 合成信号のrate--distortion。\\
10 & E5A & 小画像への凍結転移。\\
11 & E5B & 条件が整った場合のみMNIST確認。\\
\bottomrule
\end{tabularx}
\end{table}

\section{卒論最小スコープ}

最小スコープはE0--E3である。次の四点が揃えば、E4/E5が未成功でも卒論として一貫した主張を作れる。

\begin{enumerate}
  \item 一般化Boole写像の不変Cauchy尺度とLyapunov指数を再現した。
  \item TM時間特徴の可読性と識別可能性限界を正負対照で示した。
  \item 局所横安定性と大域同期成功率が一致しない条件を示した。
  \item 同期誤差と入力保持がtrade-offを持つことをN=8系で示した。
\end{enumerate}

E4/E5は「有用な圧縮」まで到達するための拡張であり、基礎仮説が通らない場合は優先度を下げる。

\section{十日間の実装例}

\begin{longtable}{p{15mm}p{37mm}p{88mm}}
\toprule
日 & 主作業 & 終了条件 \\
\midrule
Day 1 & repositoryとE0A & config、型、テスト、1条件の軌道図。\\
Day 2 & E0A本実験 & Cauchy QQ、尺度、Lyapunov、精度監査。\\
Day 3 & E0B/E1A pilot & TM直交・shift、alpha readout小実験。\\
Day 4 & E1A/E1B & 容量一致比較、rank-one下限。\\
Day 5 & E2A & clean-fit、雑音・欠測・長さ曲線。\\
Day 6 & E3A audit & 行和不変性、uniform/random pilot。\\
Day 7 & E3B & 局所指数とbasin map。\\
Day 8 & E3C & K走査、同期-情報Pareto。\\
Day 9 & E4A pilot & 三信号族、五表現、d=8/16。\\
Day 10 & レポート & 図表、失敗原因、次仮説、再現コマンド。\\
\bottomrule
\end{longtable}

\section{停止条件}

次の場合は、規模を増やす前に止まる。

\begin{itemize}
  \item E0の理論値を再現しない。
  \item rank-one負対照が理論下限を不自然に破る。
  \item 同期多様体が不変でないのに横指数を議論している。
  \item 主要結論が1 seedまたは1精度だけで成立する。
  \item decoderを大きくしたときだけ性能が出る。
  \item testを見て力学パラメータを選んでいる。
  \item clipやfallbackを外すと全結果が消える。
\end{itemize}

\section{研究仮説の更新規則}

当初の強い仮説
\begin{quote}
カオス同期すれば情報圧縮できる。
\end{quote}
は広すぎる。実験を通じて検証する更新仮説は次である。

\begin{quote}
\textbf{観測写像が識別可能性を保ち、横方向だけが有限時間で収縮し、入力が同期多様体内または遅い過渡モードへ入り、readoutが残存モードに整合するときに限り、同期縮約は復号可能な低次元表現となる。}
\end{quote}

この四条件のどれが欠けたかを、E1--E5が順に切り分ける。

\section{最終レポートの主要図}

最低限、次の六図を揃える。

\begin{enumerate}
  \item E0: 理論対数値のCauchy尺度・Lyapunov一致図。
  \item E1: TM/Fourier/shuffleのalpha readout比較。
  \item E1B: full-rankとrank-oneの識別可能性図。
  \item E3B: 局所\(\lambda_\perp\)と大域basin成功率の二段相図。
  \item E3C: 同期誤差と入力BAのK依存、またはPareto図。
  \item E4/E5: rate--distortionと復元例。未実行なら次実験計画図。
\end{enumerate}

\begin{reportbox}
図のcaptionだけで、データ、パラメータ、線・点・誤差帯の意味が分かるように書く。縦軸・横軸、単位、対数軸、seed数、集約方法を省略しない。本文では図番号で参照する。
\end{reportbox}

\begin{taskbox}
\textbf{課題 10.1}\quad 自分の実装予定を\cref{tab:implementation_order}に合わせてIssueへ分割せよ。一つのIssueは一つの反証可能な問いにする。

\medskip
\textbf{課題 10.2}\quad 最初の二週間でどこまで到達するかを、必須、余力、延期の三段階に分けて書け。
\end{taskbox}

\appendix

% ============================================================
\chapter{実装アーキテクチャとAPI契約}
\label{app:architecture}
% ============================================================

\section{推奨ディレクトリ}

\begin{Verbatim}[fontsize=\small,frame=single]
projects/chaos-sync-thesis/
|-- pyproject.toml
|-- src/chaos_sync_textbook/
|   |-- configs.py
|   |-- maps.py
|   |-- coupling.py
|   |-- features.py
|   |-- metrics.py
|   |-- datasets.py
|   `-- runners/
|       |-- e0a.py
|       |-- e1a.py
|       `-- ...
|-- configs/
|   |-- e0a.yaml
|   |-- e1a.yaml
|   `-- ...
|-- tests/
|   |-- test_maps.py
|   |-- test_features.py
|   |-- test_coupling.py
|   `-- test_metrics.py
|-- runs/
|-- reports/
`-- references/
\end{Verbatim}

\section{データクラス}

設定と結果を辞書のまま渡さず、責務ごとにdataclassへ分ける。最低限、次を持つ。

\begin{lstlisting}[style=pythonstyle,caption={設定クラスの最小例。}]
from dataclasses import dataclass
from pathlib import Path

@dataclass(frozen=True)
class OrbitConfig:
    """How: Store an immutable single-orbit simulation contract."""
    alpha: float
    burn_in: int
    n_steps: int
    seed: int
    dtype: str = "float64"

@dataclass(frozen=True)
class RunPaths:
    """How: Centralize all artifact paths for one run."""
    root: Path
    metrics_csv: Path
    summary_json: Path
    predictions_npz: Path
\end{lstlisting}

\section{局所写像API}

\begin{lstlisting}[style=pythonstyle,caption={局所写像APIの契約。}]
import numpy as np
from numpy.typing import NDArray

FloatArray = NDArray[np.float64]

def generalized_boole(x: FloatArray, alpha: float) -> FloatArray:
    """How: Apply F_alpha(x)=alpha*(x-1/x) elementwise."""
    # Why not: silently clipping x near zero changes the map itself.
    raise NotImplementedError

def boole_derivative(x: FloatArray, alpha: float) -> FloatArray:
    """How: Evaluate alpha*(1+1/x**2) elementwise."""
    raise NotImplementedError

def simulate_orbit(config: OrbitConfig, x0: float) -> FloatArray:
    """How: Simulate burn-in and return only the measured trajectory."""
    raise NotImplementedError
\end{lstlisting}

更新関数は、状態だけでなく、非有限値、特異点近傍、最大絶対値を診断情報として返す設計でもよい。例外と欠損を区別する。

\section{特徴API}

\begin{lstlisting}[style=pythonstyle,caption={TM特徴APIの契約。}]
def cayley_mode(
    x: FloatArray,
    mu: float,
    gamma: float,
    order: int,
) -> NDArray[np.complex128]:
    """How: Map a heavy-tailed orbit to the bounded TM power mode."""
    raise NotImplementedError

def complex_lag_correlation(
    z: NDArray[np.complex128],
    lags: tuple[int, ...],
    mask: NDArray[np.bool_] | None = None,
) -> NDArray[np.complex128]:
    """How: Estimate pairwise-complete complex lag correlations."""
    raise NotImplementedError

def tm_feature_vector(
    x: FloatArray,
    orders: tuple[int, ...],
    lags: tuple[int, ...],
) -> FloatArray:
    """How: Concatenate real and imaginary parts in a fixed order."""
    raise NotImplementedError
\end{lstlisting}

特徴の並び順をコード内の偶然に任せず、metadataへ保存する。例えば
\begin{center}
\code{channel=0, order=1, lag=1, part=real}
\end{center}
のようなfeature nameを生成する。

\section{結合API}

\begin{lstlisting}[style=pythonstyle,caption={row-stochastic出力混合のAPI契約。}]
def output_mixing_step(
    state: FloatArray,
    alpha: float,
    coupling: float,
    weight: FloatArray,
    common_input: float = 0.0,
) -> FloatArray:
    """How: Apply B_K F_alpha(state) plus a common scalar drive."""
    raise NotImplementedError

def assert_row_stochastic(weight: FloatArray, atol: float = 1e-12) -> None:
    """How: Validate nonnegative rows and row sums before simulation."""
    raise NotImplementedError
\end{lstlisting}

\section{runnerの責務}

各runnerは次だけを行う。

\begin{enumerate}
  \item configを読み、出力先を一意に作る。
  \item seedを設定し、データを生成する。
  \item pure function群を呼び出す。
  \item 未集約結果を保存する。
  \item 集約・図・validationを作る。
  \item configと主要artifactのhashを保存する。
\end{enumerate}

runner内へ理論式や特徴計算を直接書かない。再利用・単体テストができなくなるためである。

% ============================================================
\chapter{テスト仕様}
\label{app:tests}
% ============================================================

\section{E0テスト}

\begin{lstlisting}[style=pythonstyle,caption={理論整合性テストの例。}]
def test_boole_scale_fixed_point() -> None:
    """What: gamma_star satisfies gamma=alpha*(gamma+1/gamma)."""
    alpha = 0.5
    gamma = boole_gamma_star(alpha)
    residual = gamma - alpha * (gamma + 1.0 / gamma)
    assert abs(residual) < 1e-12


def test_cayley_mode_is_unimodular() -> None:
    """What: Every finite real input maps to the unit circle."""
    x = np.array([-10.0, -1.0, 0.0, 1.0, 10.0])
    z = cayley_mode(x, mu=0.0, gamma=1.0, order=3)
    assert np.allclose(np.abs(z), 1.0, atol=1e-12)
\end{lstlisting}

\section{同期多様体テスト}

\begin{lstlisting}[style=pythonstyle,caption={同期不変性テストの例。}]
def test_output_mixing_preserves_sync_manifold() -> None:
    """What: Identical nodes remain identical after one update."""
    n_nodes = 8
    state = np.full(n_nodes, 0.7)
    weight = np.full((n_nodes, n_nodes), 1.0 / n_nodes)
    nxt = output_mixing_step(
        state=state,
        alpha=0.5,
        coupling=0.52,
        weight=weight,
        common_input=0.1,
    )
    assert np.max(np.abs(nxt - nxt[0])) < 1e-12
\end{lstlisting}

添付原稿型のランダム加法結合には、逆に「行和が異なると同期状態が一stepで崩れる」負対照テストを作る。

\section{識別可能性テスト}

rank-one観測ではswapペアの特徴が一致することを検証する。targetやIDを渡さず、観測系列だけから特徴を作る。

\begin{lstlisting}[style=pythonstyle,caption={rank-one衝突の負対照。}]
def test_rank_one_features_collide_under_source_swap() -> None:
    """What: A symmetric one-channel observation loses source order."""
    y_ab = source_a + source_b
    y_ba = source_b + source_a
    f_ab = tm_feature_vector(y_ab, orders=(1, 2, 3, 4), lags=(1, 2))
    f_ba = tm_feature_vector(y_ba, orders=(1, 2, 3, 4), lags=(1, 2))
    assert np.allclose(f_ab, f_ba, atol=1e-12)
\end{lstlisting}

\section{漏洩テスト}

\begin{enumerate}
  \item scalerのfit indexとtest indexが交差しない。
  \item PCAのfitはtrainのみ。
  \item paired sampleは同じsplitに入る。
  \item 同じ軌道seedから切り出した窓は同じsplitに入る。
  \item test metricを使うmodel selection関数を作らない。
\end{enumerate}

\section{artifactテスト}

run終了時に次を自動検査する。

\begin{enumerate}
  \item 必須ファイルが存在する。
  \item CSVの条件数・seed数がconfigと一致する。
  \item すべての指標が有限、または失敗理由を持つ。
  \item 図の元データが保存される。
  \item hash再計算が一致する。
\end{enumerate}

% ============================================================
\chapter{主要導出の補足}
\label{app:derivations}
% ============================================================

\section{Cauchy線形安定性}

\(X_j\)の特性関数が\(\varphi_j(t)=\exp(\ii\mu_jt-\gamma_j|t|)\)なら、独立性より
\begin{align}
  \E\exp\left(\ii t\sum_ja_jX_j\right)
  &=\prod_j\varphi_j(a_jt)\\
  &=\exp\left[
      \ii t\sum_ja_j\mu_j
      -|t|\sum_j|a_j|\gamma_j
    \right].
\end{align}
これは\cref{eq:cauchy_linear_stability}のCauchy特性関数である。

\section{Boole尺度固定点}

\cref{eq:boole_scale_map}の固定点は
\begin{align}
  \gamma&=\alpha\gamma+\frac{\alpha}{\gamma},\\
  (1-\alpha)\gamma^2&=\alpha.
\end{align}
尺度は正なので正根だけを採用し、\cref{eq:boole_gamma_star}を得る。

\section{出力混合ネットワークの横指数}

同期軌道を\(\vect x_t=s_t\ones\)、微小摂動を\(\vect\delta_t\)とする。\cref{eq:output_mixing_network}を一次展開すると
\begin{equation}
  \vect\delta_{t+1}
  \approx F'_\alpha(s_t)\mat B_K\vect\delta_t.
\end{equation}
\(\mat W\vect v_q=\nu_q\vect v_q\)なら、\(\mat B_K\vect v_q=(1-K+K\nu_q)\vect v_q\)である。したがってmode振幅は\cref{eq:graph_mode_variational}に従い、対数成長率を時間平均すると\cref{eq:graph_mode_lyapunov}を得る。

\section{自己結合を除く全結合}

\begin{equation}
  \mat W_{\mathrm{no\ self}}
  =\frac{\ones\ones^\top-\mat I}{N-1}
\end{equation}
とする。同期方向では\(\nu_0=1\)、zero-sum横方向では
\begin{equation}
  \nu_\perp=-\frac1{N-1}.
\end{equation}
よって横混合倍率は
\begin{equation}
  1-K+K\nu_\perp
  =1-\frac{N}{N-1}K.
\end{equation}
下側臨界は
\begin{equation}
  K_{c,\mathrm{low}}
  =\frac{N-1}{N}\left(1-\e^{-\lambda_F}\right)
\end{equation}
となる。\(N\to\infty\)で自己結合を含む全平均混合の式へ近づく。

\section{rank-one衝突のBayes誤差}

同じ観測に\(\vect a\)と\(\vect b\)が同確率で対応する場合、二乗誤差の条件付きリスク
\begin{equation}
  R(\vect z)
  =\frac12\|\vect z-\vect a\|_2^2
   +\frac12\|\vect z-\vect b\|_2^2
\end{equation}
を\(\vect z\)で微分すると、最適解は\((\vect a+\vect b)/2\)である。\(\vect a=(a,b)\)、\(\vect b=(b,a)\)を代入すれば、各座標MSEは\((a-b)^2/4\)となる。

% ============================================================
\chapter{主張監査表}
\label{app:audit}
% ============================================================

\begin{longtable}{p{44mm}p{27mm}p{72mm}}
\caption{本研究での主張状態。}\label{tab:claim_audit}\\
\toprule
主張 & 状態 & 本研究での扱い \\
\midrule
\endfirsthead
\toprule
主張 & 状態 & 本研究での扱い \\
\midrule
\endhead
Cauchy特性関数と独立線形和の安定性 & \Known & E0の数学的基礎として使用。\\
Cauchyから単位円へのCayley変換とTM直交性 & \Known & 紙上導出とE0B数値検証。\\
一般化Boole写像のCauchy尺度写像 & \SourceClaim & E0Aで独立軌道から数値再現。\\
Boole写像のLyapunov閉形式 & \SourceClaim & E0Aで時間平均と照合。\\
無限次元ランダム結合の一変数尺度縮約 & \SourceClaim & 独立性、quenched disorder、row sum、有限\(N\)を監査。\\
臨界値が\(\E|\varepsilon|\)だけで決まる & \SourceClaim & E3Aで同じ絶対一次moment・異なる符号分布を比較。\\
K-fold cotangentの\(UM_k=M_{Kk}\) & \SourceClaim/基準系 & E0Bの正対照として使う。\\
PF prime-factor spectrum、KMS、mass gap & \SourceClaim & 卒論の中核実験では仮定しない。\\
AKOrNが同期を学習表現へ利用できる & 先行研究の実験結果 & 動的block・stimulus・readout設計の参考。カオス性能の証拠にはしない。\\
output-mixingのgraph mode横指数 & \Derived & E3Cで数値照合。\\
rank-one観測のswap誤差下限 & \Derived & E1Bの負対照。\\
臨界近傍で同期と情報保持が両立する & \Hypothesis & E3C以降で検証。\\
TMが全baselineより優れる & \Hypothesis & E1/E2/E4/E5で条件別に検証し、普遍主張を避ける。\\
\bottomrule
\end{longtable}

% ============================================================
\chapter{実験ノートとレポートのテンプレート}
\label{app:templates}
% ============================================================

\section{実験前}

\begin{enumerate}
  \item EIDと一文の研究質問。
  \item 帰無仮説と対立仮説。
  \item 操作変数、固定変数、乱数単位。
  \item train/validation/testの生成規則。
  \item 主要指標、副指標、失敗率。
  \item \Gate と停止条件。
  \item 予定図表。
  \item 計算budget。
\end{enumerate}

\section{実験後}

\begin{enumerate}
  \item 実行日時、commit hash、環境。
  \item configからの変更。変更がなければ「なし」。
  \item 未集約結果ファイルへの参照。
  \item 主要数値と95\%区間。
  \item \Gate ごとの支持・未支持。
  \item 理論との一致・不一致。
  \item 数値誤差、漏洩、選択バイアスの監査。
  \item 次に検証する単一仮説。
\end{enumerate}

\section{レポート本文の骨格}

\begin{Verbatim}[fontsize=\small,frame=single]
1. はじめに
   1.1 背景
   1.2 研究質問と仮説
2. 原理
   2.1 局所写像と不変測度
   2.2 同期多様体と横Lyapunov指数
   2.3 TM読み出しと識別可能性
3. 実験方法
   3.1 データ生成と分割
   3.2 実装・パラメータ・seed
   3.3 評価指標と事前Gate
4. 結果
   4.1 正対照・負対照
   4.2 主要結果
   4.3 感度分析
5. 考察
   5.1 理論との比較
   5.2 代替説明
   5.3 限界
6. 結論
7. 参考文献
\end{Verbatim}

\section{図captionの例}

\begin{quote}
図X.X: \(\alpha=1/2\)、\(N=8\)の出力混合Boole系における結合強度\(K\)と入力保持性能。左は終端24stepの同期RMS中央値、右はgroup split testにおけるTM--graph logistic probeのbalanced accuracyである。点は10 seed平均、帯はseed bootstrap 95\%区間を表す。縦破線は本モデルの局所臨界値\(K_c=0.5\)である。
\end{quote}

% ============================================================
\chapter{記号表}
\label{app:symbols}
% ============================================================

\begin{longtable}{p{28mm}p{113mm}}
\toprule
記号 & 意味 \\
\midrule
\(\alpha\) & 一般化Boole写像の局所パラメータ。\\
\(\beta\) & 生tangent写像\(\tan(\beta x)\)のパラメータ。\\
\(K\) & 結合強度。モデルごとに意味が異なるためモデル名を併記する。\\
\(N\) & ノード数。\\
\(\mu,\gamma\) & Cauchy分布の位置・尺度。\\
\(\gamma_*\) & 尺度写像の正固定点。\\
\(F_\alpha\) & 一般化Boole写像。\\
\(T_\beta\) & 生tangent写像。\\
\(H_K\) & K-fold cotangent写像。\\
\(q_{\mu,\gamma}\) & Cayley変換。\\
\(M_k\) & TM power basisの第\(k\)モード。\\
\(\lambda_F\) & 単一Boole写像のLyapunov指数。\\
\(\lambda_\parallel\) & 同期多様体に沿うLyapunov指数。\\
\(\lambda_\perp\) & 同期多様体に横断するLyapunov指数。\\
\(E_{\mathrm{sync}}\) & ノード間RMS同期誤差。\\
\(\nu_q\) & 結合行列\(\mat W\)の第\(q\)固有値。\\
\(d_z\) & 潜在次元。\\
\(\NMSE{}\) & 分散正規化平均二乗誤差。\\
\(\effrank{}\) & 特異値participation ratioによるeffective rank。\\
\bottomrule
\end{longtable}

% ============================================================
\backmatter
\chapter*{参考文献}
\addcontentsline{toc}{chapter}{参考文献}
\begin{thebibliography}{99}

\bibitem{mathlab2024}
京都大学 工学部情報学科 数理工学コース 実験演習ワーキンググループ編,
『数理工学実験 2024年度版』, 2024.

\bibitem{infinite_sync}
K. Umeno,
``Infinite Dimensional Chaotic Synchronization for Randomly Coupling Systems and Dynamical Phase Transition,''
manuscript dated July 5, 2026.

\bibitem{mass_gap}
K. Umeno,
``Universal Mass Gap Formula $\Delta m=m_0\ln K$: Spontaneous Kubo--Martin--Schwinger State Emergence with Maximal Quantum Chaos,''
manuscript dated July 30, 2026.

\bibitem{akorn}
T. Miyato, S. L\"owe, A. Geiger, and M. Welling,
``Artificial Kuramoto Oscillatory Neurons,''
International Conference on Learning Representations, 2025.

\bibitem{umeno1998}
K. Umeno,
``Superposition of chaotic processes with convergence to L\'evy's stable law,''
Physical Review E, Vol. 58, p. 2644, 1998.

\bibitem{umeno_okubo2016}
K. Umeno and K.-i. Okubo,
``Exact Lyapunov exponents of the generalized Boole transformations,''
Progress of Theoretical and Experimental Physics, 2016, 021A01, 2016.

\bibitem{kano_umeno2022}
T. Kano and K. Umeno,
``Chaotic synchronization of mutually coupled systems with arbitrary proportional linear relations,''
Chaos, Vol. 32, 113137, 2022.

\bibitem{kaneko1990}
K. Kaneko,
``Clustering, coding, switching, hierarchical ordering, and control in a network of chaotic elements,''
Physica D, Vol. 41, pp. 137--172, 1990.

\bibitem{kuramoto1984}
Y. Kuramoto,
Chemical Oscillations, Waves, and Turbulence,
Springer, 1984.

\bibitem{wood2006}
木下是雄,
『理科系の作文技術』,
中央公論新社, 1981.

\end{thebibliography}

\chapter*{最後に}
\addcontentsline{toc}{chapter}{最後に}

本研究で最も重要なのは、同期、カオス、圧縮、復号を一つの魅力的な言葉へまとめないことである。各概念を別の式と別の指標で測り、その間の論理を実験でつなぐ。理論値に合わない結果、TMがbaselineに負ける結果、深い同期で情報が消える結果も、原因を切り分けて保存すれば研究成果である。

最初の実装は\EID{E0A-BOOLE-LOCAL-VALIDATION}である。\cref{eq:boole_map,eq:boole_gamma_star,eq:empirical_boole_lyap}を自分で実装し、Cauchy QQ、MAD、Lyapunov一致図の三つを出すところから始める。

\end{document}
